\graphicspath{{notes/2inner/asy/}}
\def\at#1#2{#1_{#2}}

\setcounter{section}{1}

\section{Inner Product Spaces, part 1}

You should be familiar with the \emph{scalar/dot product} in $\R^2$. For any vectors $\vx=\stwovec{x_1}{x_2},\ \vy=\stwovec{y_1}{y_2}$, define\smallbreak
\begin{minipage}[t]{0.65\linewidth}\vspace{-15pt}
\[\vx\cdot\vy:=x_1y_1+x_2y_2\]
The \emph{norm} or \emph{length} of a vector is $\Nm{\vx}=\sqrt{\vx\cdot\vx}$.\smallbreak
The \emph{angle} $\theta$ between vectors satisfies $\vx\cdot\vy=\Nm\vx\Nm\vy\cos\theta$.\smallbreak
Vectors are \emph{orthogonal} or \emph{perpendicular} precisely when $\vx\cdot\vy=0$.
\end{minipage}\begin{minipage}[t]{0.34\linewidth}\vspace{0pt}
\flushright\includegraphics{dot-angle}
\end{minipage}\smallbreak
The dot product is what allows us to compute lengths of and angles between vectors in $\R^2$. An \emph{inner product} is an algebraic structure that generalizes this idea to other vector spaces.

\subsection{Real and Complex Inner Product Spaces}\label{sec:realcomplexips}

Unless explicitly stated otherwise, throughout this chapter $\F$ is either the real $\R$ or complex $\C$ field.

\begin{defn}{}{innerproduct}
An \emph{inner product space} $(V,\ip{\ ,\ })$ is a vector space $V$ over $\F$ together with an \emph{inner product}: a function $\ip{\ ,\ }:V\times V\to\F$ satisfying the following properties $\forall\vx,\vy,\vz\in V,\ \lambda\in\F$,
\begin{enumeratea}\itemsep2pt
  \item \emph{Linear}:\quad $\ip{\lambda\vx+\vy,\vz}=\lambda\ip{\vx,\vz}+\ip{\vy,\vz}$
  \item \emph{Conjugate-Symmetric}:\quad $\ip{\vy,\vx}=\cl{\ip{\vx,\vy}}$ \hfill (complex conjugate!)
  \item \emph{Positive-definite}:\quad $\vx\neq \V0\implies \ip{\vx,\vx}>0$
\end{enumeratea}
The \emph{norm} or \emph{length} of $\vx\in V$ is $\Nm{\vx}:=\sqrt{\ip{\vx,\vx}}$. A \emph{unit vector} has $\Nm\vx=1$.\smallbreak
Vectors $\vx,\vy$ are \emph{perpendicular/orthogonal} if $\ip{\vx,\vy}=0$ and \emph{orthonormal} if they are additionally unit vectors.
\end{defn}



\boldsubsubsection{Real inner product spaces}

The definition simplifies slightly when $\F=\R$.
\begin{itemize}
  \item Conjugate-symmetry becomes plain \emph{symmetry}: $\ip{\vy,\vx}=\ip{\vx,\vy}$.
  \item Linearity $+$ symmetry yields \emph{bilinearity}: a real inner product is also linear in its second slot
  \[\ip{\vx,\lambda\vy+\vz}=\lambda\ip{\vx,\vy}+\ip{\vx,\vz}\]
\end{itemize}
A real inner product is often termed a \emph{positive-definite, symmetric, bilinear form.}\medbreak

The simplest example is the natural generalization of the dot product.

\begin{defn}{}{}
\emph{Euclidean space} means $\R^n$ equipped with the \emph{standard inner (dot) product},
\[\ip{\vx,\vy} =\vx\cdot\vy=\vy^T\vx=\sum_{j=1}^nx_jy_j=x_1y_1+\cdots+x_ny_n,\qquad \Nm{\vx}=\sqrt{\sum_{j=1}^nx_j^2}\]
\end{defn}

Unless the inner product is stated explicitly, if we refer to $\R^n$ as an inner product space, we mean Euclidean space. However, there are many other ways to make $\R^n$ into an inner product space\ldots
\goodbreak

\begin{example}{}{ipsweight}
Define an alternative inner product on $\R^2$ via
	\[\ip{\vx,\vy}=x_1y_1+3x_2y_2\]
	It is easy to check that this satisfies the required properties: 
	\begin{enumeratea}
  	\item Linearity: follows from the associative/distributive laws in $\R$,
  	\begin{align*}
  		\ip{\lambda\vx+\vy,\vz}&=(\lambda x_1+y_1)z_1+3(\lambda x_2+y_2)z_2
  		=\lambda(x_1z_1+3x_2z_2)+(y_1z_1+3y_2z_2)\\
  		&=\lambda\ip{\vx,\vz}+\ip{\vy,\vz}
  	\end{align*}
  	\item Symmetry: follows from the commutativity of multiplication in $\R$,
  	\[\ip{\vy,\vx}=y_1x_1+3y_2x_2=x_1y_1+3x_2y_2=\ip{\vx,\vy}\]
		\item Positive-definiteness:\quad if $\vx\neq\V0$, then $\ip{\vx,\vx}=x_1^2+3x_2^2>0$.
	\end{enumeratea}
	With respect to $\ip{\ ,\ }$, the concept of \emph{orthogonality} feels strange: e.g., $\{\vx,\vy\}=\left\{\frac 12\stwovec 11,\frac 1{2\sqrt{3}}\stwovec 3{-1}\right\}$ is an orthonormal set!
  \[
  	\Nm\vx^2=\frac 14(1^2+3\cdot 1^2)=1,\quad \Nm\vy^2=\frac 1{12}(3^2+3\cdot 1^2)=1,\quad \ip{\vx,\vy}=\frac 1{4\sqrt 3}(3-3)=0
  \]
	However, with respect to the standard dot product, these vectors are not special:
  \[\vx\cdot\vx=\frac 12,\qquad \vy\cdot\vy=\frac 56,\qquad \vx\cdot \vy=\frac 1{2\sqrt 3}\]
  We have the same \emph{vector space} $\R^2$, but different \emph{inner product spaces}: $(\R^2,\ip{\ ,\ })\neq (\R^2,\cdot)$.
\end{example}

The above is an example of a \emph{weighted inner product}: choose \emph{weights} $a_1,\ldots,a_n\in\R^+$ and define
	\[\ip{\vx,\vy}=\sum_{j=1}^na_jx_jy_j=a_1x_1y_1+\cdots+a_nx_ny_n\]
	It is a simple exercise to check that this defines an inner product on $\R^n$. In particular, $\R^n$ may be equipped with \emph{infinitely many distinct inner products!}\medbreak

More generally, a symmetric matrix $A\in M_n(\R)$ is \emph{positive-definite} if $\vx^TA\vx>0$ for all non-zero $\vx\in\R^n$. It is straightforward to check that
	\[\ip{\vx,\vy}:=\vy^TA\vx\]
	defines an inner product on $\R^n$. In fact \emph{all} inner products on $\R^n$ arise in this fashion! The weighted inner products correspond to $A$ being diagonal (Euclidean space is $A=I$), but this is not required.
	
\begin{example}{}{ipsweight3}
	The matrix $A=\begin{smatrix}
	3&1\\1&1
	\end{smatrix}$ is positive-definite and thus defines an inner product
	\[\ip{\vx,\vy}=3x_1y_1+x_1y_2+x_2y_1+x_2y_2\]
\end{example}

\goodbreak


\begin{lemm}{Basic properties}{ipbasic}
Let $V$ be an inner product space, let $\vx,\vy,\vz\in V$ and $\lambda\in\F$.\smallbreak

  
  \begin{minipage}[b]{0.7\linewidth}\vspace{0pt}
  \begin{enumerate}\itemsep2pt
  \item $\ip{\V0,\vx}=0$
  \item $\Nm\vx=0\iff \vx=\V0$
  \item $\Nm{\lambda\vx}=\nm{\lambda}\Nm{\vx}$
  \item $\ip{\vx,\vz}=\ip{\vy,\vz}$ for all $\vz\implies \vx=\vy$
  	\item (Cauchy--Schwarz inequality)\quad $\nm{\ip{\vx,\vy}}\le\Nm\vx\Nm\vy$ with equality if and only if $\vx,\vy$ are parallel
    \item (Triangle Inequality)\quad $\Nm{\vx+\vy}\le\Nm\vx+\Nm\vy$ with equality if and only if
    $\vx,\vy$ are parallel and point in the same direction
\end{enumerate}
  \end{minipage}
  \hfill
  \begin{tabular}[b]{@{}c@{}}
  	\includegraphics{dot-triangle}\\
  	Triangle inequality
  \end{tabular}
\end{lemm}

Be careful with notation: $\nm\lambda$ means the \emph{absolute value/modulus} (of a \emph{scalar}), while $\Nm{\vx}$ means the \emph{norm} (of a \emph{vector}).\smallbreak
In the real case, Cauchy--Schwarz allows us to define \emph{angle} via $\cos\theta=\frac{\ip{\vx,\vy}}{\Nm\vx\Nm\vy}$, since the right-hand side lies in the interval $[-1,1]$. However, except in Euclidean $\R^2$ and $\R^3$, this notion is of limited use; \emph{orthogonality} ($\ip{\vx,\vy}=0$) and \emph{orthonormality} are usually all we care about.

\begin{proof}
Parts 1--3 are exercises. For simplicity, we prove 5 and 6 only when $\F=\R$.
\begin{enumerate}\setcounter{enumi}{3}
%   \item $\ip{\V0,\vx}=\ip{\vx-\vx,\vx}=\ip{\vx,\vx}-\ip{\vx,\vx}=0$
%   \item ($\Rightarrow$) is the contrapositive of defining property (c), while $(\Leftarrow)$ follows from part 1.
%   \item $\Nm{\lambda\vx}^2=\ip{\lambda\vx,\lambda\vx}=\lambda\cl\lambda\ip{\vx,\vx}=\nm{\lambda}^2\Nm{\vx}^2$
  \item Let $\vz=\vx-\vy$, apply the linearity condition and part 2:
  \[\ip{\vx,\vz}=\ip{\vy,\vz}\implies 0=\ip{\vx-\vy,\vz}=\ip{\vx-\vy,\vx-\vy}=\Nm{\vx-\vy}^2\implies \vx=\vy\]
  \item If $\vy=\V0$, the result is trivial. WLOG (and by part 3) we may assume $\Nm\vy=1$; if the inequality holds for this, then it holds for all non-zero $\vy$ by parts 1 and 4. Now expand:
  \begin{align*}
  0&\le\Nm{\vx-\ip{\vx,\vy}\vy}^2\tag{positive-definiteness}\\
  &=\Nm\vx^2+\nm{\ip{\vx,\vy}}^2\Nm{\vy}^2 -\ip{\vx,\vy}\ip{\vx,\vy} -\ip{\vx,\vy}\ip{\vy,\vx} \tag{bilinearity}\\
  &=\Nm\vx^2-\nm{\ip{\vx,\vy}}^2\tag{symmetry}
  \end{align*}
  Taking square-roots establishes the inequality.\smallbreak
  By part 2, equality holds if and only if $\vx=\ip{\vx,\vy}\vy$, which is precisely when $\vx,\vy$ are linearly dependent.
  \item We establish the squared result.
  \begin{align*}
  \bigl(\Nm\vx+\Nm\vy\bigr)^2-\Nm{\vx+\vy}^2&=\Nm\vx^2+2\Nm\vx\Nm\vy+\Nm\vy^2-\Nm\vx^2-\ip{\vx,\vy}-\ip{\vy,\vx}-\Nm\vy^2\\
  &=2\bigl(\Nm\vx\Nm\vy-\ip{\vx,\vy}\bigr)\\
  &\ge 2\bigl(\Nm\vx\Nm\vy-\nm{\ip{\vx,\vy}}\bigr)\\
  &\ge 0 \tag{Cauchy--Schwarz}
  \end{align*}
  Equality requires both equality in Cauchy--Schwarz ($\vx,\vy$ parallel) and that $\ip{\vx,\vy}\ge 0$; since $\vx,\vy$ are already parallel, this means that one is a non-negative multiple of the other.\qedhere
\end{enumerate}
\end{proof}

\boldsubsubsection{Complex Inner Product Spaces}

Definition \ref{defn:innerproduct} is already set up nicely when $\C=\F$. One subtle difference comes from how we expand linear combinations in the second slot.

\begin{lemm}{}{}
An inner product is a \emph{positive-definite, conjugate-symmetric, \textcolor{red}{sesquilinear}\footnotemark form}: it is \emph{conjugate-linear} (\emph{anti-linear}) in the second slot,
\[\ip{\vz,\lambda\vx+\vy}=\textcolor{red}{\cl\lambda}\ip{\vz,\vx}+\ip{\vz,\vy}\]
\end{lemm}

The proof is very easy if you remember your complex conjugates; try it!\smallbreak

\footnotetext{The prefix \emph{sesqui-} means \emph{one-and-a-half}; for instance a \emph{sesquicentenary} is a 150 year anniversary.} 

\textcolor{red}{Warning!} If you dabble in the dark arts of Physics, be aware that their convention\footnote{The common Physics notation relates to ours via $\ip{\vx\,|\,\vy}=\ip{\vy,\vx}$.} is for an inner product to be conjugate-linear in the \emph{first} entry and linear in the second!

\begin{defn}{}{}
The \emph{standard (Hermitian) inner product} and \emph{norm} on $\C^n$ are
\[\ip{\vx,\vy}=\vy^*\vx=\sum_{j=1}^nx_j\cl{y_j}=x_1\cl{y_1}+\cdots+x_n\cl{y_n},\qquad \Nm{\vx}=\sqrt{\sum_{j=1}^n\nm{x_j}^2}\]
where $\vy^*=\cl\vy^T$ is the \emph{conjugate-transpose} of $\vy$ and $\nm{x_j}$ is the \emph{modulus.}
\end{defn}

Weighted inner products may be defined as on page \pageref{ex:ipsweight}:
\[\ip{\vx,\vy}=\sum_{j=1}^na_jx_j\cl{y_j}=a_1x_1\cl{y_1}+\cdots+a_nx_n\cl{y_n}\]
Note that the weights $a_j$ must still be \emph{positive real numbers.} We may similarly define inner products in terms of positive-definite matrices $\ip{\vx,\vy}=\vy^*A\vx$.

\begin{defn}{}{}
A matrix $A\in M_n(\C)$ is \emph{Hermitian (self-adjoint)} if $A^*=A$. It is moreover \emph{positive-definite} if $\vx^*A\vx>0$ for all non-zero $\vx\in\C^n$.
\end{defn}

The self-adjoint condition reduces to symmetry ($A^T=A$) if $A$ is a real matrix.

\begin{example}{}{}
It can be seen (Exercise \ref{exs:complexposdef}) that $A=\begin{smatrix}
3&-i\\i&3
\end{smatrix}$ is positive-definite, whence
\[\ip{\vx,\vy}=\bigl(\cl{y_1}\ \ \cl{y_2}\bigr)
\begin{pmatrix}
3&-i\\i&3
\end{pmatrix}
\twovec{x_1}{x_2} =
3x_1\cl{y_1}+ix_1\cl{y_2}-ix_2\cl{y_1}+3x_2\cl{y_2}\]
defines an inner product on $\C^2$.
\end{example}

Almost all results in this chapter will be written for general inner product spaces, thus covering the real and complex cases simultaneously. If you don't feel confident with complex numbers, simply let $\F=\R$ and delete all complex conjugates at first read! Very occasionally a different proof will be required depending on the field. For simplicity, examples will more often use real inner products. 

\goodbreak

\boldsubsubsection{Further Examples} As before, the field $\F$ must be either $\R$ or $\C$.

\begin{defn}{Frobenius inner product}{}
If $A,B\in M_{m\times n}(\F)$, define
\[\ip{A,B}=\tr(B^*A)\]
where $\tr$ is the \emph{trace} of an $n\times n$ matrix; this makes $M_{m\times n}(\F)$ into an inner product space.
\end{defn}

This isn't really a new example: if we map $A\to\F^{m\times n}$ by stacking the columns of $A$, then the Frobenius inner product is the standard inner product in disguise.


\begin{example}{}{frobenius}
In $M_{3\times 2}(\C)$,
\begin{gather*}
\begin{aligned}
\ip{\begin{pmatrix}
1&i\\
2-i&0\\
0&-1
\end{pmatrix}
,
\begin{pmatrix}
0&7\\
1&2i\\
3-2i&4
\end{pmatrix}}
&=\tr\begin{pmatrix}
0&1&3+2i\\
7&-2i&4
\end{pmatrix}\begin{pmatrix}
1&i\\
2-i&0\\
0&-1
\end{pmatrix}
=\tr\begin{pmatrix}
2-i&-3-2i\\
9-4i&-4+7i
\end{pmatrix}\\
&=2-i-4+7i=-2-8i
\end{aligned}\\
\begin{aligned}
\Nm{\begin{pmatrix}
1&i\\
2-i&0\\
0&-1
\end{pmatrix}}^2
&=\tr\begin{pmatrix}
1&2+i&0\\
-i&0&-1
\end{pmatrix}\begin{pmatrix}
1&i\\
2-i&0\\
0&-1
\end{pmatrix}
=\tr\begin{pmatrix}
6&i\\
-i&2
\end{pmatrix}=8
\end{aligned}
\end{gather*}
\end{example}


\begin{defn}{$L^2$ inner product}{}
Given a \emph{real} interval $[a,b]$, the function
\[\ip{f,g}:=\int_a^b f(t)\cl{g(t)}\,\dt\]
defines an inner product on the space $C[a,b]$ of continuous functions $f:[a,b]\to\F$.
\end{defn}

With careful restriction, this works even for infinite intervals and a larger class of functions.\footnote{For us, functions will always be continuous (often polynomials) on closed bounded intervals. The \emph{\href{https://en.wikipedia.org/wiki/Square-integrable_function}{square-integrable} functions} and \emph{$L^2$-spaces} for which the inner product is named are a more complicated business and beyond this course.} Verifying the required properties is straightforward if you know a little analysis; for instance continuity allows us to conclude
\[\Nm f^2=\int_a^b\nm{f(x)}^2\,\dx=0\iff f(x)\equiv 0\]
This is our first example of an \emph{infinite-dimensional} inner product space.

\begin{example}{}{}
Let $f(x)=x$ and $g(x)=x^2$; these lie in the inner product space $C[-1,1]$ with respect to the $L^2$ inner product.
\[\ip{f,g}=\int_{-1}^1x^3\,\dx=0,\qquad \Nm f^2=\int_{-1}^1x^2\,\dx=\frac 23,\qquad \Nm g^2=\int_{-1}^1x^4\,\dx=\frac 25\]
With some simple scaling, we see that $\left\{\frac 1{\Nm f}f,\frac 1{\Nm g}g\right\}=\left\{\sqrt{\frac 32}x,\sqrt{\frac 52}x^2\right\}$ forms an orthonormal set.
\end{example}
\goodbreak

\begin{defn}{$\ell^2$ inner product}{ell2ip}
A sequence $(x_n)$ is \emph{square-summable} if \smash[b]{$\sum\limits_{n=1}^\infty\nm{x_n}^2<\infty$.} These sequences form a vector space on which we can define an inner product\footnotemark
\[\ip{(x_n),(y_n)}=\sum_{n=1}^\infty x_n\cl{y_n}\]
\end{defn}

\footnotetext{Neither of these facts are obvious; in particular, we'd need to see that the sum of two square-summable sequences is also square-summable.}

In essence we've taken the standard inner product on $\F^n$ and let $n\to\infty$! This example, and its $L^2$ cousin, are the prototypical \emph{\href{https://en.wikipedia.org/wiki/Hilbert_space}{Hilbert spaces,}} which have great application to differential equations, signal processing, etc. Since a rigorous discussion requires a significant amount of analysis (convergence of series, completeness, integrability), these objects are generally beyond the course.\medbreak

Our final example of an inner product is a useful, and hopefully obvious, hack to which we shall repeatedly appeal in examples.

\begin{lemm}{}{basisonip}
Let $V$ be a vector space over $\R$ or $\C$. If $\beta$ is a basis of $V$, then there exists exactly one inner product on $V$ for which $\beta$ is an orthonormal set.
\end{lemm}



\begin{exercises}
\exstart Evaluate the inner product of the given vectors.
\begin{enumerate}\setcounter{enumi}{1}
  \item[]\begin{enumerate}
    \item $\vx=\sthreevec 121$, $\vy=\sthreevec{1}{-1}1$ where $\ip{\vx,\vy}=2x_1y_1+3x_2y_2+x_3y_3$
    \item $\vx=\stwovec 1{2i}$, $\vy=\stwovec{5i}{4}$ where $\ip{\vx,\vy}$ is the standard Hermitian inner product on $\C^2$
    \item $\vx=\stwovec 1{2i}$, $\vy=\stwovec{5i}{4}$ where $\ip{\vx,\vy}=\vy^*\begin{smatrix}
    2&-i\\i&2
    \end{smatrix}\vx$
    \item $f(x)=x-1$, $g(x)=x+1$ where $\ip{f,g}$ is the $L^2$ inner product on $C[0,2]$
  \end{enumerate}
  
  \item Suppose $\ip{\vx,\vy}=\ip{\vx,\vz}$ for all $\vx$. Prove that $\vy=\vz$.
  
  \item For each $\vz\in V$, The linearity condition says that the map $\rT_\vz:V\to\F$ defined by
  \[\rT_\vz(\vx)=\ip{\vx,\vz}\]
  is linear ($\rT_\vz$ is an element of the \emph{dual space} $V^*$). What, if anything, can you say about the function $\rU_\vz:\vx\mapsto\ip{\vz,\vx}$?
  
  \item Define $\ip{\vx,\vy}:=\sum\limits_{j=1}^nx_jy_j$ on $\C^n$. Is this an inner product? Which of the properties (a), (b), (c) from Definition \ref{defn:innerproduct} does it satisfy?
  
  \item\begin{enumerate}
    \item Verify that the matrix in Example \ref{ex:ipsweight3} is positive-definite and that  $\ip{\vx,\vy}=3x_1y_1+x_1y_2+x_2y_1+x_2y_2$ therefore defines an inner product on $\R^2$.\\
  (\emph{Hint: Try to write $\Nm{\vx}^2$ as a sum of squares\ldots})
  	\item Let $\vx=\stwovec 11$. With respect to the inner product in part (a), find a non-zero unit vector $\vy$ which is \emph{orthogonal} to $\vx$.
  \end{enumerate}
  \goodbreak
  
  \item\label{exs:complexposdef} By multiplying out $\nm{x_1-ix_2}^2$, show that the matrix $\begin{smatrix}3&-i\\i&3\end{smatrix}$ is positive-definite.\smallbreak
  (\emph{Hint: recall that $\nm{a}^2=a\cl a$ for complex numbers!!!})
  \goodbreak
  
  \item Show that every eigenvalue of a positive definite matrix is positive.
  
  \item Prove parts 1, 2 and 3 of Lemma \ref{lemm:ipbasic}.
  
  \item\label{ex:pythag} Let $V$ be an inner product space, prove Pythagoras' Theorem:
  \[\text{If }\ip{\vx,\vy}=0,\text{ then }\Nm{\vx+\vy}^2=\Nm{\vx}^2+\Nm{\vy}^2\]
  
	\item Use basic algebra to prove the Cauchy--Schwarz inequality for vectors $\vx=\stwovec{a}{b}$ and $\vy=\stwovec{c}{d}$ in $\R^2$ with the standard (dot) product.
	
	\item Prove the Cauchy--Schwarz and triangle inequalities for a \emph{complex} inner product space. What has to change compared to the proof of Lemma \ref{lemm:ipbasic}?
	
	\item Prove the polarization identities:
	\begin{enumerate}
	  \item In any real inner product space: $\ip{\vx,\vy}=\frac 14\Nm{\vx+\vy}^2 -\frac 14\Nm{\vx-\vy}^2$
	  \item In any complex inner product space: $\ip{\vx,\vy}=\frac 14\sum\limits_{k=1}^4i^k\Nm{\vx+i^k\vy}^2$
	\end{enumerate}
	\emph{If you know the length of every vector, then you know the inner product!}
	
	\item Prove that $\int_0^2\frac{\sqrt x}{x+1}\,\dx\le \frac 2{\sqrt 3}$\smallbreak
	(\emph{Hint: use Cauchy--Schwarz})

  
  \item\label{ex:fourier} Let $m\in\Z$ and consider the complex-valued function $f_m(x)=\frac 1{\sqrt{2\pi}}e^{i mx}$. If $\ip{\ ,\ }$ is the $L^2$ inner product on $C[-\pi,\pi]$, prove that $\{f_m:m\in\Z\}$ is an orthonormal set.\smallbreak
  \emph{This example is central to the study of \emph{Fourier series.}}
  
  (\emph{Hint: If complex functions are scary, use Euler's formula $e^{i mx}=\cos mx+i\sin mx$ and work with the real-valued functions $\cos mx$ and $\sin mx$. The difficulty is that you then need integration by parts\ldots})
  
  \item Let $\ip{\ ,\ }$ be an inner product on $\F^n$ (recall that $\F=\R$ or $\C$). Define the matrix $A\in M_n(\F)$ by $A_{jk}=\ip{\ve_k,\ve_j}$ where $\{\vect en\}$ is the standard basis. Verify that $A$ is the \emph{matrix of the inner product}:
  \[\forall \vx,\vy\in\F^n,\ \ip{\vx,\vy}=\vy^*A\vx\]
	In particular,
  \begin{itemize}
    \item $A$ is a \emph{Hermitian/self-adjoint} matrix: $A^*=A$ (if $\F=\R$ this is simply \emph{symmetric});
    \item $A$ is \emph{positive-definite}: for all $\vx\in\F^n$, $\vx^*A\vx>\V0$;
  \end{itemize}
  \emph{More generally, if $\beta=\{\vect vn\}$ is a basis then $A_{jk}=\ip{\vv_k,\vv_j}$ defines the matrix of the inner product with respect to $\beta$:
  \[\ip{\vx,\vy}=[\vy]_\beta^*A[\vx]_\beta\]}

\end{enumerate}
\end{exercises}
\clearpage


\subsection{Orthogonal Sets and the Gram--Schmidt Process}\label{sec:gs}

% 
% 
% The Gram--Schmidt process is a simple algorithm whose purpose is the construction of \emph{orthogonal} sets of vectors. It relies on a basic property of orthogonal sets.

We start with a simple definition, relating subspaces to orthogonality. 

\begin{defn}{}{}
Let $U$ be a subspace of an inner product space $V$. The \emph{orthogonal complement} $U^\perp$ is the set
\[U^\perp=\{\vx\in V:\forall\vu\in U,\ \ip{\vx,\vu}=0\}\]
\end{defn}

It is easy to check that $U^\perp$ is itself a subspace of $V$ and that $U\cap U^\perp=\{\V0\}$. It can moreover be seen that $U\subseteq (U^\perp)^\perp$, though equality need not hold in infinite dimensions (see Exercise \ref{ex:ell2ex}).

\begin{example}{}{orthcompeeasy}
$U=\Span\left\{\sthreevec 100,\sthreevec 0{-1}3\right\}\le\R^3$ has orthogonal complement $U^\perp=\Span\left\{\sthreevec 031\right\}$.
\end{example}

\begin{minipage}[t]{0.62\linewidth}\vspace{0pt}
As in the example, we often have a direct sum decomposition $V=U\oplus U^\perp$: otherwise said,
\[\forall \vx\in V, \exists \text{ unique }\vu\in U,\vw\in U^\perp \text{ such that }\vx=\vu+\vw\]
In such a case, we'll call $\vu=\pi_U(\vx)$ and $\vw=\pi_U^\perp(\vx)$ the \emph{orthogonal projections} of $\vx$ onto $U$ and $U^\perp$ respectively. As our first result shows, these are easy to compute whenever $U$ has a finite orthogonal basis.
\end{minipage}
\begin{minipage}[t]{0.37\linewidth}\vspace{0pt}
\flushright%\includemedia[transparent=false, activate=onclick, add3Djscript=asylabels.js, %add3Djscript=3Dspintool.js,
% 3Dmenu,
%   3Dcoo=6.2 46.8 5.9,
%   3Dc2c=0.51 -0.77 0.39,
%   3Droo=325,
%   3Droll=-6.2,
% 3Dlights=Headlamp]{\includegraphics{gram-orthproj3+0_0}}{gram-orthproj3+0.prc}
\href{https://www.math.uci.edu/~ndonalds/math121b/gram-orthproj.html}{\includegraphics{gram-orthproj}}
\end{minipage}

\begin{thm}{}{orthcalc}
Let $V$ be an inner product space and let $U=\Span \beta$ where $\beta=\{\vect{\vu}{n}\}$ is an orthogonal set of non-zero vectors;
%\[\ip{\vu_j,\vu_k}=0\text{ if $j\neq k$}\]
\[\ip{\vu_j,\vu_k}=\begin{cases}
0&\text{if $j\neq k$}\\
\Nm{\vu_j}^2\neq 0&\text{if $j=k$}
\end{cases}\]
Then:
\begin{enumerate}
  \item $\beta$ is a basis of $U$ and each $\vx\in U$ has unique representation
  \[\vx=\sum\limits_{j=1}^n\frac{\ip{\vx,\vu_j}}{\Nm{\vu_j}^2}\vu_j\tag{$\ast$}\]
  This simplifies to $\vx=\sum\ip{\vx,\vu_j}\vu_j$ if $\beta$ is an orthonormal set.
  \item $V=U\oplus U^\perp$. For any $\vx\in V$, we may write $\vx=\vu+\vw$ where
  \[\vu=\sum\limits_{j=1}^n\frac{\ip{\vx,\vu_j}}{\Nm{\vu_j}^2}\vu_j\in U\quad\text{and}\quad \vw=\vx-\vu\in U^\perp\]
  %are the orthogonal projections.
\end{enumerate}
\end{thm}


Observe that ($\ast$) essentially calculates the co-ordinate vector $[\vx]_\beta\in\F^n$. Recalling how unpleasant such calculations have been in the past, often requiring large matrix inversions, we immediately see the power of inner products and orthogonal bases. 
\goodbreak
 

\begin{proof}
\begin{enumerate}
  \item Since $\beta$ spans $U$, a given $\vx\in U$ may be written
  \[\vx=\sum_{k=1}^na_k\vu_k\]
  for some scalars $a_k$. The orthogonality of $\beta$ recovers the required expression for $a_j$:
  \[\ip{\vx,\vu_j}=\sum_{k=1}^na_k\ip{\vu_k,\vu_j}=a_j\Nm{\vu_j}^2\]
 	Finally, let $\vx=\V0$ to see that $\beta$ is linearly independent.
  \item Clearly $\vu\in U$. For each $\vu_k$, the orthogonality of $\beta$ tells us that
  \[\ip{\vw,\vu_k}=\ip{\vx-\vu,\vu_k}=\ip{\vx,\vu_k}-\sum_{j=1}^n\frac{\ip{\vx,\vu_j}}{\Nm{\vu_j}^2}\ip{\vu_j,\vu_k}=\ip{\vx,\vu_k}-\ip{\vx,\vu_k}=0\]
  Since $\vw$ is orthogonal to a basis of $U$ it is orthogonal to any element of $U$; we conclude that $\vw\in U^\perp$. Finally, $U\cap U^\perp=\{\V0\}$ forces the uniqueness of the decomposition $\vx=\vu+\vw$, whence $V=U\oplus U^\perp$.\qedhere
\end{enumerate}
\end{proof}


\begin{examples}{}{}
\exstart Consider the standard orthonormal basis $\beta=\{\ve_1,\ve_2\}$ of $\R^2$. For any $\vx=\stwovec{x_1}{x_2}$, we easily check that
  \[\sum_{j=1}^2\ip{\vx,\ve_j}\ve_j=x_1\ve_1+x_2\ve_2=\vx\]
\begin{enumerate}\setcounter{enumi}{1}
  \item In $\R^3$, $\beta=\{\vu_1,\vu_2,\vu_3\}=\left\{\sthreevec 123,\sthreevec 2{-1}0,\sthreevec 36{-5}\right\}$ is an orthogonal set and thus a basis. We compute the co-ordinates of $\vx=\sthreevec 742$ with respect to $\beta$:\vspace{-3pt}
	\begin{align*}
	\vx&=\sum_{j=1}^3\frac{\ip{\vx,\vu_j}}{\Nm{\vu_j}^2}\vu_j = \frac{7+8+6}{1+4+9}\threevec 123 + \frac{14-4+0}{4+1+0}\threevec 2{-1}0 + \frac{21+24-10}{9+36+25}\threevec 36{-5}\\
	&= \frac 32\threevec 123 + 2\threevec 2{-1}0 + \frac 12\threevec 36{-5} \implies [\vx]_\beta=\threevec{3/2}{2}{1/2}
	\end{align*}
	Compare this with the painfully slow augmented matrix method for finding co-ordinates!
% 	
% 	\item $S=\{1,x,x^2-\frac 13\}$ is orthogonal with respect to the $L^2$ inner product on $[-1,1]$. Moreover,
% 	\[\Nm{1}^2=\int_{-1}^11^2\,\dx=2,\qquad \Nm{x}^2=\int_{-1}^1x^2\,\dx=\frac 23,\qquad \Nm{x^2-\frac 13}^2=\frac 8{45}\]
% 	We confirm that
% 	\[x^2+x=\]

	\item Revisiting Example \ref{ex:orthcompeeasy}, let $\vx=\sthreevec 111$. Since $\beta=\left\{\sthreevec 100,\sthreevec 0{-1}3\right\}$ is an orthogonal basis of $U$, we observe that
	\[\pi_U(\vx)=\frac{1}{1}\threevec 100+\frac{2}{10}\threevec 0{-1}3=\frac 15\threevec 5{-1}3,\qquad \pi_U^\perp(\vx)=\vx-\pi_U(\vx)=\frac 25\threevec 031\]
	are the orthogonal projections corresponding to $\R^3=U\oplus U^\perp$.
\end{enumerate}
\end{examples}

\boldsubsubsection{The Gram--Schmidt Process}

Theorem \ref{thm:orthcalc} tells us how to compute the orthogonal projections corresponding to $V=U\oplus U^\perp$, \emph{provided $U$ has a finite, orthogonal basis.} Given how useful such bases are, our next goal is to see that such exist for any finite-dimensional subspace. Helpfully there exists a constructive algorithm.

\begin{thm}{Gram--Schmidt}{gs}
Suppose $S=\{\vect{\vs}{n}\}$ is a linearly independent subset of an inner product space $V$. Construct a sequence of vectors $\vect{\vu}{n}$ inductively:
\begin{itemize}
  \item Choose $\vu_1=a_1\vs_1$ where $a_1\neq 0$\vspace{-2pt}
  \item For each $k\ge 2$, choose $\displaystyle\vu_k=a_k\left(\vs_k-\sum_{j=1}^{k-1}\frac{\ip{\vs_k,\vu_j}}{\Nm{\vu_j}^2}\vu_j\right)$ where $a_k\neq 0$ \hfill($\dag$)
\end{itemize}
Then $\beta:=\{\vect{\vu}{n}\}$ is an orthogonal basis of $\Span S$.
\end{thm}

The purpose of the scalars $a_k$ is to give you some freedom; choose them to avoid unpleasant fractions! If you want a set of \emph{orthonormal} vectors, it is easier to scale everything after the algorithm is complete. Indeed, by taking $S$ to be a basis of $V$ and normalizing the resulting $\beta$, we conclude:

\begin{cor}{}{orthbasis}
Every finite-dimensional inner product space has an orthonormal basis.
\end{cor}

\begin{example}{}{gs1}
$S=\{\vs_1,\vs_2,\vs_3\}=\left\{\sthreevec 100,\sthreevec 2{-1}3,\sthreevec 111\right\}$ is a linearly independent subset of $\R^3$.
\begin{enumerate}
  \item Choose $\vu_1=\vs_1=\sthreevec 100\implies \Nm{\vu_1}^2=1$
  \item $\vs_2-\frac{\ip{\vs_2,\vu_1}}{\Nm{\vu_1}^2}\vu_1=\sthreevec 2{-1}3-\frac{2}{1}\sthreevec 100 =\sthreevec 0{-1}3:\text{ choose }\vu_2=\sthreevec 0{-1}3\implies \Nm{\vu_2}^2=10$
  \item $\vs_3-\frac{\ip{\vs_3,\vu_1}}{\Nm{\vu_1}^2}\vu_1 -\frac{\ip{\vs_3,\vu_2}}{\Nm{\vu_2}^2}\vu_2=\sthreevec 111-\frac{1}{1}\sthreevec 100 -\frac{2}{10}\sthreevec 0{-1}3 =\frac 25\sthreevec 0{3}{1}:\text{ choose }\vu_3=\sthreevec 03{1}$
\end{enumerate}
% 	\begin{gather*}
% \vu_1=\vs_1=\sthreevec 100\implies \Nm{\vu_1}^2=1\\
% \vs_2-\frac{\ip{\vs_2,\vu_1}}{\Nm{\vu_1}^2}\vu_1=\sthreevec 2{-1}3-\frac{2}{1}\sthreevec 100 =\sthreevec 0{-1}3:\text{ choose }\vu_2=\sthreevec 0{-1}3\implies \Nm{\vu_2}^2=10\\
% \vs_3-\frac{\ip{\vs_3,\vu_1}}{\Nm{\vu_1}^2}\vu_1 -\frac{\ip{\vs_3,\vu_2}}{\Nm{\vu_2}^2}\vu_2=\sthreevec 111-\frac{1}{1}\sthreevec 100 -\frac{2}{10}\sthreevec 0{-1}3 =\frac 25\sthreevec 0{3}{1}:\text{ choose }\vu_3=\sthreevec 03{1}
% 	\end{gather*}
The orthogonality of $\beta=\{\vu_1,\vu_2,\vu_3\}$ is clear. It is now trivial to observe that $\left\{\vu_1,\frac 1{\sqrt{10}}\vu_2,\frac 1{\sqrt{10}}\vu_3\right\}$ is an orthonormal basis of $\R^3$.
\end{example}


\begin{proof}[Proof of Theorem \ref{thm:gs}]
For each $k\le n$, define $S_k=\{\vect{\vs}{k}\}$ and $\beta_k=\{\vect{\vu}{k}\}$. We prove by induction that each $\beta_k$  is an orthogonal set of non-zero vectors and that $\Span\beta_k=\Span S_k$. The Theorem is then the terminal case $k=n$.\smallbreak
(Base case $k=1$)\quad Certainly $\beta_1=\{a_1\vs_1\}$ is an orthogonal set and $\Span\beta_1=\Span S_1$.\\[5pt]
(Induction step)\quad Fix $k\ge 2$, assume $\beta_{k-1}$ is an orthogonal non-zero set and that $\Span \beta_{k-1}=\Span S_{k-1}$. By Theorem \ref{thm:orthcalc}, $\vu_k\in(\Span\beta_{k-1})^\perp$. We also see that $\vu_k\neq \V0$, for if not,
\[(\dag)\implies\vs_k\in\Span \beta_{k-1}=\Span S_{k-1}\]
and $S$ would be linearly dependent. It follows that $\beta_k$ is an orthogonal set of non-zero vectors.\smallbreak
Moreover, $\vs_k\in\Span \beta_k\implies \Span S_k\le\Span \beta_k$. Since these spaces have the same (finite) dimension $k$, we conclude that $\Span \beta_k=\Span S_k$.\smallbreak
By induction, $\beta$ is an orthogonal, non-zero spanning set for $\Span S$; by Theorem \ref{thm:orthcalc}, it is a basis.
\end{proof}
\goodbreak

\begin{example}{}{ell201}
This time we work in the space of polynomials $P(\R)$ equipped with the $L^2$ inner product $\ip{f,g}=\int_0^1f(x)g(x)\,\dx$ on the interval $[0,1]$. Let $S=\{1,x,x^2\}$ and apply the algorithm:
\begin{enumerate}
	\item Choose $f_1(x)=1 \implies \Nm{f_1}^2=\int_0^11\,\dx=1$
	\item $x-\frac{\ip{x,f_1}}{\Nm{f_1}^2}f_1=x-\int_0^1x\,\dx =x-\frac 12$\smallbreak
	We choose $f_2(x)=2x-1$, with $\Nm{f_2}^2=\int_0^1\left(2x-1\right)^2\,\dx=\frac 1{3}$
	\item	$x^2-\frac{\ip{x^2,f_1}}{\Nm{f_1}^2}f_1-\frac{\ip{x^2,f_2}}{\Nm{f_2}^2}=
	x^2-\int_0^1x^2\,\dx-\frac{\int_0^1x^2\left(2x-1\right)\,\dx}{1/3}\left(2x-1\right)=x^2-x+\frac 16$\smallbreak
	We choose $f_3(x)=6x^2-6x+1$ with $\Nm{f_3}^2=\int_0^1\left(6x^2-6x+1\right)^2\,\dx=\frac 1{5}$
	\end{enumerate}
	It follows that $\Span S$ has an orthonormal basis
	\[\beta=\left\{1,\sqrt 3(2x-1),\sqrt 5(6x^2-6x+1)\right\}\]
	This example can be extended to arbitrary degree since the countable set $\{1,x,x^2,\ldots\}$ is basis of $P(\R)$. Indeed this shows that $(P(\R),\ip{\ ,\ })$ has an orthonormal basis. 
\end{example}









Gram--Schmidt also shows that our earlier discussion of orthogonal projections is generic.
  
\begin{cor}{}{}
If $U$ is a finite-dimensional subspace of an inner product space $V$, then $V=U\oplus U^\perp$ and the orthogonal projections may be computed as in Theorem \ref{thm:orthcalc}.
\end{cor}

If $U$ is an infinite-dimensional subspace of $V$, then we need not have $V=U\oplus U^\perp$ and the orthogonal projections might not be well-defined (see, for example, Exercises \ref{ex:ell2ex} and \ref{ex:fourierx}). Instead, if $\beta$ is an orthonormal basis of $U$, it is common to describe the coefficients $\ip{\vx,\vu}$ for each $\vu\in\beta$ as the \emph{Fourier coefficients}, and the infinite sum
\[\sum_{\vu\in\beta}\ip{\vx,\vu}\vu\]
as the \emph{Fourier series} of $\vx$, provided the sum converges.

% We finally check that everything goes through in a general inner product space.
% 
% \begin{thm}{}{projcleanup}
% The functions $\pi_U$ and $\pi^\perp_U$ are linear maps satisfying
% \[\pi_U(\vx)=\begin{cases}
%   \vx&\text{if }\vx\in U\\
%   \V0&\text{if }\vx\in U^\perp
%   \end{cases}\qquad \pi_U^\perp(\vx)=\begin{cases}
%   \vx&\text{if }\vx\in U^\perp\\
%   \V0&\text{if }\vx\in U
%   \end{cases}\]
%   Moreover $\pi_U^\perp(\vw) \in U^\perp$ for any $\vw\in V$ and $U\oplus U^\perp=V$.\\[5pt]
%   Finally, $\pi_U(\vx)$ is the \emph{closest} vector in the subspace $U$ to $\vx$, in that
%   \[\forall\vw\in U,\ \Nm{\pi_U(\vx)-\vx}\le\Nm{\vw-\vx}\]
% \end{thm}
% 
% \begin{proof}
% $\pi_U$ is linear by property (a) of Definition \ref{defn:innerproduct}. The linearity of $\pi_U^\perp=\I-\pi_U$ is then immediate.\\[5pt]
% The values for $\pi_U(\vx)$ and $\pi^\perp_U(\vx)$ follow trivially from Definition \ref{defn:orthproj} and ($\ast$) in Theorem \ref{thm:orthproj}.\\[5pt]
% Finally, let $\vw\in V$, and let $U=\Span\{\vect{\vw}{n}\}$ be spanned by orthonormal vectors: then we have a calculation very like part of the proof of Gram--Schmidt:
% \begin{align*}
% \pi_U^\perp(\vw)=\vw-\sum_{j=1}^n\ip{\vw,\vw_j}\vw_j \implies \ip{\vw_i,\pi_U^\perp(\vw)} &= \ip{\vw_i,\vw}-\sum_{j=1}^n\cl{\ip{\vw,\vw_j}}\ip{\vw_i,\vw_j}\\
% &=\ip{\vw_i,\vw}-\cl{\ip{\vw,\vw_i}}\ip{\vw_i,\vw_i} =0
% \end{align*}
% Since $\pi_U^\perp(\vw)$ is orthogonal to a basis of $U$, it must lie in $U^\perp$.\\[5pt]
% We can now decompose any vector into $U$- and $U^\perp$-parts: $\vx=\pi_U(\vx)+\pi_U^\perp(\vx)$. Since $U\cap U^\perp=\{\V0\}$, we see that $U\oplus U^\perp=V$.\\[5pt]
% Finally, use Pythagoras' (Exercise \ref{ex:pythag}) to observe that
% \[\Nm{\vw-\vx}^2=\Nm{\vw-\pi_U(\vx)-\pi_U^\perp(\vx)}^2=\Nm{\vw-\pi_U(\vx)}^2+\Nm{\pi_U^\perp(\vx)}^2\]
% whence
% \[\Nm{\vw-\vx}\ge \Nm{\vw-\pi_U(\vx)}\tag*{\qedhere}\]
% \end{proof}
% 
% 
% \emph{Warning}: the Theorem \emph{only} applies when $\dim U<\infty$ since the proof depends crucially on being able to define $\pi_U$ for which we require a \emph{finite} orthogonal basis. Indeed there are examples where $V\neq U\oplus U^\perp$.\\[-5pt]
% 
% The final property is often known as the \emph{least-squares} property of projections since one is minimizing distance: for the standard inner product on $\R^n$, this amounts to minimizing the distance squared, which is itself a sum of squares. The method of least-squares curve-fitting in statistics can be rephrased as a linear algebra problem involving projections.
  
\bigskip

\begin{exercises}{}{}
\exstart Apply Gram--Schmidt to obtain an orthogonal basis $\beta$ for $\Span S$. Then obtain the co-ordinate representation (Fourier coefficients) of the given vector with respect to $\beta$.
\begin{enumerate}\setcounter{enumi}{1}
  \item[]\begin{enumerate}
    \item $S=\left\{\sthreevec 111,\sthreevec 011,\sthreevec 001\right\}\subseteq\R^3$,\quad $\vx=\sthreevec 101$
    \item $S=\left\{
    \begin{smatrix}
    3&5\\-1&1
    \end{smatrix},
    \begin{smatrix}
    -1&9\\5&-1
    \end{smatrix},
    \begin{smatrix}
    7&-17\\2&-6
    \end{smatrix}
    \right\}\subseteq M_2(\R)$,\quad $X=\begin{smatrix}
    -1&27\\-4&8
    \end{smatrix}$\quad  (use the \hyperref[ex:frobenius]{Frobenius} product)
    \item $S=\left\{\sthreevec 1i0,\sthreevec 01i,\sthreevec 001\right\}\subseteq\C^3$ with $\vx=\sthreevec 111$
    \item $S=\{1,x,x^2\}$ with $\ip{f,g}=\int_{-1}^1f(x)g(x)\,\dx$ and $f(x)=x^2$.
  \end{enumerate}
  \emph{Important!} You'll likely need \emph{much} more practice than this to get comfortable with Gram--Schmidt; make up your own problems!
  \goodbreak
  
  \item Let $S=\{\vs_1,\vs_2\}=\left\{\sthreevec 103,\sthreevec 2{-1}0\right\}$ and $U=\Span S\le\R^3$. Find $\pi_U(\vx)$ if $\vx=\sthreevec 3{-1}{-2}$.\smallbreak
  (\emph{Hint: First apply Gram--Schmidt})
 	
 	\item Find the orthogonal complement to $U=\Span\{x^2\}\le P_2(\R)$ with respect to the inner product $\ip{f,g}=\int_0^1 f(t)g(t)\,\dt$
    
  \item\label{ex:fdipsmatrix} Let $\rT\in\cL(V,W)$ where $V,W$ are inner product spaces with orthonormal bases $\beta=\{\vect{\vv}{n}\}$ and $\gamma=\{\vect wm\}$ respectively. Prove that the matrix $A=[\rT]_\beta^\gamma \in M_{m\times n}(\F)$ of $\rT$ with respect to these bases has $jk^\text{th}$ entry
  \[A_{jk}=\ip{\rT(\vv_k),\vw_j}\]
  
  \item Suppose that $\beta$ is an orthonormal basis of an $n$-dimensional inner product space $V$. Prove that,
	\[\forall\vx,\vy\in V,\ \ip{\vx,\vy}=[\vy]_\beta^*[\vx]_\beta\]
	\emph{Otherwise said, the co-ordinate isomorphism $\phi_\beta:V\to \F^n$ defined by $\phi_\beta(\vx)=[\vx]_\beta$ is an isomorphism of inner product spaces where we use the standard inner product on $\F^n$}
 	
  \item\label{ex:uperpperp} Let $U$ be a subspace of an inner product space $V$. Prove the following:
  \begin{enumerate}
    \item $U^\perp$ is a subspace of $V$.
    \item $U\cap U^\perp=\{\V0\}$
    \item $U\subseteq (U^\perp)^\perp$
    \item If $V=U\oplus U^\perp$, then $U=(U^\perp)^\perp$ (this is always the case when $\dim U<\infty$)
  \end{enumerate}
  
  
  
  \item\label{ex:ell2ex} Let $\ell^2$ be the set of square-summable sequences of real numbers (Definition \ref{defn:ell2ip}). Consider the sequences $\vu_1,\vu_2,\vu_3,\ldots$, where $\vu_j$ is the zero sequence except for a single 1 in the $j^\text{th}$ entry. For instance,
 	\[\vu_4=(0,0,0,1,0,0,0,0,\ldots)\]
  \begin{enumerate}
    \item Let $U=\Span\{\vu_j:j\in\N\}$. Prove that $U^\perp$ contains only the zero sequence.
    \item Show that the sequence $\vy=\left(\frac 1n\right)$ lies in $\ell^2$, but does not lie in $U$.\smallbreak
    \emph{$U$ is therefore a \emph{proper subset} of $(U^\perp)^\perp=\ell^2$ and $\ell^2\neq U\oplus U^\perp$.}
  \end{enumerate}
 
  \item\label{ex:fourierx} Recall Exercise \hyperref[ex:fourier]{\ref*{sec:realcomplexips}.\ref*{ex:fourier}} where we saw that the set $\beta=\{\frac 1{\sqrt{2\pi}}e^{i mx}:m\in\Z\}$ is orthonormal with respect to $\ip{f,g}=\int_{-\pi}^\pi f(t)\cl{g(t)}\,\dt$.
  \begin{enumerate}
    \item Show that the \emph{Fourier series} of $f(x)=x$ is
  	\[\mathcal{F}(x):=\sum_{m=-\infty}^\infty \ip{x,\tfrac 1{\sqrt{2\pi}}e^{i mx}}\tfrac 1{\sqrt{2\pi}}e^{i mx} %=\frac 1{2\pi} \sum_{m=0}^\infty\left(\int_{-\pi}^\pi xe^{-imx}\,\dx\right)e^{imx}
  	=\sum_{m=1}^\infty\frac{2(-1)^{m+1}}m\sin mx\]
  	\item Briefly explain why the Fourier series is \emph{not} an element of $\Span\beta$.
  	\item Sketch a few of the Fourier approximations (sum up to $m=5$ or 7\ldots) and observe, when extended to $\R$, how they approximate a discontinuous periodic function.
 	\end{enumerate}

  
  \item (Hard)\quad Much of Theorem \ref{thm:orthcalc} remains true, with suitable modifications, even if $\beta$ is an \emph{infinite} set. Restate and prove as much as you can, and identify the false part(s).%\smallbreak
  %(\emph{Hint: a linear combination is always a finite sum!})

  
  % 	\item\label{ex:Tinj} Let $T\in\cL(V)$ be a linear operator on an inner product space and suppose that $\Nm{T(\vx)}=\Nm{\vx}$ for all $\vx\in V$. Prove that $T$ is injective (1--1).	
	
% 	\item Let $A,B\in M_n(\F)$ and let $\ip{\ ,\ }$ be the standard inner product on $\F^n$.
% 	\begin{enumerate}
% 	  \item Prove that $\forall\vx,\vy\in \F^n$, we have $\ip{\vx,A\vy}=\ip{A^*\vx,\vy}$.
% 	  \item Suppose $\ip{\vx,A\vy}=\ip{B\vx,\vy}$ for all $\vx,\vy\in\F^n$: prove that $B=A^*$.
% 	  \item Let $\beta$ be any orthonormal basis of $\F^n$ and let $Q$ be the matrix whose columns are the co-ordinates of the elements of $\beta$ with respect to the standard basis. Prove that $Q^*=Q^{-1}$.
% 	\end{enumerate}
\end{enumerate}
\end{exercises}

\clearpage


\subsection{The Adjoint of a Linear Operator}

Recall how the standard inner product on $\F^n$ may be written in terms of the conjugate-transpose%\footnote{Remember you can delete the complex conjugate and work in $\R^n$ at a first read. The example should help!}
\[\ip{\vx,\vy}=\vy^*\vx =\cl{\vy}^T\vx\]
We start by inserting a matrix into this expression and interpreting in two different ways. Suppose $A\in M_{m\times n}(\F)$, $\vv\in\F^n$ and $\vw\in\F^m$, then
\[\underbrace{\ip{A^*\vw,\vv}}_{\text{in }\F^n}=\vv^*(A^*\vw)= (\vv^*A^*)\vw =(A\vv)^*\vw=\underbrace{\ip{\vw,A\vv}}_{\text{in }\F^m} \tag{$\dag$}\]

\begin{example}{}{adjointeasy}
As a sanity check, let $A=\begin{smatrix}
1&2\\0&3
\end{smatrix}\in M_2(\R)$, $\vw=\stwovec xy$ and $\vv=\stwovec pq$. Then,
\begin{gather*}
\ip{A^T\twovec xy,\twovec pq}
=\ip{\begin{pmatrix}
1&0\\2&3
\end{pmatrix}\twovec xy,\twovec pq}
=\ip{\twovec x{2x+3y},\twovec pq}
=xp+(2x+3y)q\\[5pt]
\ip{\twovec xy,A\twovec pq}
=\ip{\twovec xy,\begin{pmatrix}
1&2\\0&3
\end{pmatrix}\twovec pq}
=\ip{\twovec xy,\twovec{p+2q}{3q}}
=x(p+2q)+3yq
\end{gather*}
\end{example}

Note how the inner products are evaluated on \emph{different spaces.} At the level of linear maps this is a relationship between $\rL_A\in\cL(\F^n,\F^m)$ and $\rL_{A^*}\in\cL(\F^m,\F^n)$, one that is easily generalizable.

\begin{defn}{}{adjoint}
Let $\rT\in\cL(V,W)$ where $V,W$ are inner product spaces over the same field $\F$. The \emph{adjoint} of $\rT$ is a function $\rT^*:W\to V$ (read `$\rT$-star') satisfying
\[\forall \vv\in V,\vw\in W,\quad \ip{\rT^*(\vw),\vv}=\ip{\vw,\rT(\vv)}\]
Note that the first inner product is computed within $V$ and the second within $W$.
\end{defn}

The adjoint effectively extends the conjugate-transpose to linear maps. We now use the same notation for three objects, so be careful!
\begin{itemize}
  \item If $A$ is a real or complex matrix, then $A^*=\cl A^T$ is its \emph{conjugate-transpose.}
  \item If $\rT$ is a linear map, then $\rT^*$ is its \emph{adjoint.}
  \item If $V$ is a vector space, then $V^*=\cL(V,\F)$ is its \emph{dual space.}
\end{itemize}
Thankfully the two notations line up nicely, as part 3 of our first result shows.

\begin{thm}{Basic Properties}{matrixip}
\exstart If an adjoint exists,\footnotemark then it is unique and linear.
\begin{enumerate}\setcounter{enumi}{1}
  \item If $\rT$ and $\rS$ have adjoints, then
  \[(\rT^*)^*=\rT,\quad (\rT\rS)^*=\rS^*\rT^*,\quad(\lambda\rT+\rS)^*=\cl\lambda\rT^*+\rS^*\]
  \item Suppose $V,W$ are finite-dimensional with orthonormal bases $\beta,\gamma$ respectively. Then the matrix of the adjoint of $\rT\in\cL(V,W)$ is the conjugate-transpose of the original: $[\rT^*]_\gamma^\beta=([\rT]_\beta^\gamma)^*$.
\end{enumerate}
\end{thm}

\footnotetext{Existence of adjoints is trickier, so we postpone this a little: see Corollary \ref{cor:adjexist} and Exercise \ref{exs:adjnonexist}.}
\goodbreak

\begin{proof}
\begin{enumerate}
  \item (\emph{Uniqueness})\quad Suppose $\rT^*$ and $\rS^*$ are adjoints of $\rT$. Then
	\[\ip{\rT^*(\vx),\vy}=\ip{\vx,\rT(\vy)}=\ip{\rS^*(\vx),\vy}\]
	Since this holds for all $\vy$, Lemma \ref{lemm:ipbasic} part 4 says that $\forall\vx,\ \rT^*(\vx)=\rS^*(\vx)$, whence $\rT^*=\rS^*$.\smallbreak
	(\emph{Linearity})\quad Simply translate across, use the linearity of $\rT$, and again appeal to Lemma \ref{lemm:ipbasic}:
	\begin{gather*}
		\begin{aligned}
		\forall\vz,\ \ip{\rT^*(\lambda\vx+\vy),\vz} &=\ip{\lambda\vx+\vy,\rT(\vz)} =\lambda\ip{\vx,\rT(\vz)} +\ip{\vy,\rT(\vz)} =\lambda\ip{\rT^*(\vx),\vz} +\ip{\rT^*(\vy),\vz}\\
		&=\ip{\lambda \rT^*(\vx)+\rT^*(\vy),\vz}
		\end{aligned}\\
	\implies \rT^*(\lambda\vx+\vy)=\lambda \rT^*(\vx)+\rT^*(\vy)
	\end{gather*}

	\item These may be proved similarly to part 1 and are left as an exercise.
  
  \item By Exercise \hyperref[ex:fdipsmatrix]{\ref*{sec:gs}.\ref*{ex:fdipsmatrix}}, the $jk^\text{th}$ entry of $[\rT^*]_\gamma^\beta$ is 
	\[\ip{\rT^*(\vw_k),\vv_j}=\ip{\vw_k,\rT(\vv_j)}=\cl{\ip{\rT(\vv_j),\vw_k}}=\cl{A_{kj}}\tag*{\qedhere}\]
\end{enumerate}
\end{proof}

We revisit our motivating set-up ($\dag$) in the language of part 3. Suppose:
\begin{itemize}
  \item $V=\F^n$ and $W=\F^m$ have standard orthonormal bases $\beta=\{\vect en\}$ and $\gamma=\{\vect em\}$.
  \item $\rT=\rL_A\in\cL(\F^n,\F^m)$. 
\end{itemize} Since the matrix of $\rT$ with respect to the standard bases is simply $A$ itself, the theorem confirms our earlier observation that the adjoint of $\rL_A$ is left multiplication by the conjugate-transpose $A^*$:
\[[\rT^*]_\gamma^\beta=([\rT]_\beta^\gamma)^*=A^*=[\rL_{A^*}]_\gamma^\beta\implies \rT^*=(\rL_A)^*=\rL_{A^*}\]

Here is another straightforward example, this time using the standard inner products on $\C^2$ and $\C^3$.
 
\begin{example}{}{}
Let $\rT=\rL_A\in\cL(\C^3,\C^2)$ where $A=\begin{smatrix}
  -i&1&-3\\
  2&1-i&4+2i
\end{smatrix}$.\smallbreak
Plainly $A=[\rT]_\beta^\gamma$ with respect to $\beta=\{\ve_1,\ve_2,\ve_3\}$ and $\gamma=\{\ve_1,\ve_2\}$. We conclude that $\rT^*=\rL_{A^*}$:
  \[[\rT^*]_\gamma^\beta=A^*=\begin{pmatrix}
  i&2\\
  1&1+i\\
  -3&4-2i
  \end{pmatrix}\]
  As a sanity check, multiply out a few examples of $\ip{A^*\vw,\vv}=\ip{\vw,A\vv}$; make sure you're comfortable with the fact that the left inner product is on $\C^2$ and the right on $\C^3$!
\end{example}

The Theorem tells us that every linear map $\rT\in\cL(V,W)$ between finite-dimensional spaces has an adjoint and moreover how to compute it:
\begin{enumerate} 
  \item Choose orthonormal bases (exist by Corollary \ref{cor:orthbasis}) and find the matrix $[\rT]_\beta^\gamma$.
  \item Take the conjugate-transpose $\bigl([\rT]_\beta^\gamma\bigr)^*$ and translate back to find $\rT^*\in\cL(W,V)$.
\end{enumerate}
The prospect of twice applying Gram--Schmidt and translating between linear maps and their matrices is unappealing; calculating this way can quickly become an enormous mess! In practice, it is often better to try a modified approach; see for instance part 2(b) of the next Example.
\goodbreak



\begin{examples}{}{deriv2d}
Let $\rT=\diff x\in\cL(P_1(\R))$ be the derivative operator; $\rT(a+bx)=b$. We treat $P_1(\R)$ as an inner product space in two ways.
\begin{enumerate}
  \item Equip the inner product for which the standard basis $\epsilon=\{1,x\}$ is orthonormal. Then
  \[[\rT]_\epsilon=
  \begin{pmatrix}
  0&1\\0&0
  \end{pmatrix}
  \implies [\rT^*]_\epsilon=
  \begin{pmatrix}
  0&0\\1&0
  \end{pmatrix}
  \implies \rT^*(a+bx)=ax\]
  
  \item Equip the $L^2$ inner product $\ip{f,g}=\int_0^1f(x)g(x)\,\dx$. As we saw in Example \ref{ex:ell201}, the basis $\beta=\{f_1,f_2\}=\{1,2x-1\}$ is orthogonal with $\Nm{f_1}=1$ and $\Nm{f_2}=\frac 1{\sqrt 3}$. We compute the adjoint of $\rT=\diff x$ in two different ways.
  \begin{enumerate}
    \item The basis $\gamma=\{g_1,g_2\}=\left\{f_1,\sqrt 3f_2\right\}=\left\{1,\sqrt 3(2x-1)\right\}$ is orthonormal. Observe that
    \begin{gather*}
    \rT(g_1)=0,\quad \rT(g_2)=2\sqrt 3\implies [\rT]_\gamma=\begin{pmatrix}
    0&2\sqrt 3\\
    0&0
    \end{pmatrix}\implies [\rT^*]_\gamma=\begin{pmatrix}
    0&0\\
    2\sqrt 3&0
    \end{pmatrix}\\
    \begin{aligned}
    \implies \rT^*(a+bx)&=\rT^*\left(a+\frac b2+\frac b{2\sqrt 3}\sqrt 3(2x-1)\right) =\left(a+\frac b2\right)\cdot 2\sqrt 3g_2\\
    &=3(2a+b)(2x-1)
    \end{aligned}
    \end{gather*}
    \item Use the orthogonal basis $\beta$ and the projection formula (Theorem \ref{thm:orthcalc}). With $p(x)=a+bx$,
		\begin{align*}
		\rT^*(p)&=\frac{\ip{\rT^*(p),f_1}}{\Nm{f_1}^2}f_1 +\frac{\ip{\rT^*(p),f_2}}{\Nm{f_2}^2}f_2 =\ip{p,\rT(1)} +\ip{p,\rT(2x-1)}\cdot 3(2x-1)\\
		&=\ip{p,0} +3\ip{p,2}(2x-1) =3\left(\int_0^12(a+bx)\,\dx\right)(2x-1) \\
		&=3(2a+b)(2x-1)
		\end{align*}
		Note the advantages here: no square roots and no need to change basis at the end!
  \end{enumerate}
\end{enumerate}
  The calculations for the second example were much nastier, even though we were already in possession of an orthogonal basis. The crucial point is that the two examples produce \emph{different maps} $\rT^*$: the adjoint depends on the inner product! 
\end{examples}




\boldsubsubsection{Why should we care about adjoints?}

Adjoints might seem merely to be an abstraction of something simple (transposes) for its own sake. A convincing explanation of why adjoints are useful takes a lot of work; here is a short version.\smallbreak

Given a linear map $\rT\in\cL(V)$ on an inner product space, we now have \emph{two desirable types of basis.}
\begin{enumerate}
  \item Eigenbasis: diagonalizes $\rT$.
  \item Orthonormal basis: recall (Theorem \ref{thm:orthcalc}) how these simplify computations.
\end{enumerate}
The capstone result of this course is the famous \emph{spectral theorem} (Theorem \ref{thm:spec}) which says, in short, that \emph{\href{https://en.wikipedia.org/wiki/self-adjoint_operator}{self-adjoint operators}} ($\rT^*=\rT$) have an \emph{orthonormal eigenbasis,} the holy grail of easy computation! Such operators are important both theoretically and in applications such as quantum mechanics.
\goodbreak



\boldsubsubsection{The Fundamental Subspaces Theorem}

To every linear map are associated its range and nullspace. These interact nicely with the adjoint\ldots

\begin{thm}{}{fundssthm}
If $\rT\in\cL(V,W)$ has adjoint $\rT^*$, then,
  \begin{enumerate}
    \item $\cR(\rT^*)^\perp=\cN(\rT)$
    \item If $V$ is finite dimensional, then $\cR(\rT^*)=\cN(\rT)^\perp$
  \end{enumerate}
  The corresponding results hold if we swap $V\leftrightarrow W$ and $\rT\leftrightarrow \rT^*$.
\end{thm}

The proof is left to Exercise \ref{ex:fundssthm}. You've likely observed this with transposes of small matrices.

\begin{example}{}{}
Let $A=\begin{smatrix}
1&2&-1\\
0&3&-2
\end{smatrix}$. Viewed as a linear map between Euclidean spaces, $\rT=\rL_A$ has adjoint $\rT^*=\rL_{A^T}$. It is easy to compute the relevant subspaces:
\begin{gather*}
\cR(A)=\R^2,\quad \cN(A^T)=\{\V0\},\qquad \cR(A^T)=\Span\left\{\sthreevec 12{-1},\sthreevec 03{-2}\right\},\quad \cN(A)=\Span\sthreevec{-1}23
\end{gather*}
\end{example}


% \begin{proof}
% $\vv\in\cR(\rT^*)^\perp\iff \forall\vw\in W,\ \ip{\vv,\rT^*(\vw)}=0 \iff \forall\vw\in W,\ \ip{\rT(\vv),\vw}=0 \iff \rT(\vv)=\V0$
% \end{proof}


\boldsubsubsection{The Riesz Representation Theorem}

This powerful result demonstrates a natural relation between an inner product space and its \emph{dual-space} $V^*=\cL(V,\F)$.

%When $V$ is finite-dimensional, we've seen how orthonormal bases convert calculations with adjoints to matrices, though this is often ugly and obscures the forest for the trees. To help remedy this, we observe a natural relation between $V$ and its \emph{dual-space} $V^*=\cL(V,\F)$.

\begin{thm}{}{riesz}
If $V$ is finite-dimensional and $g:V\to\F$ is linear, then there exists a unique $\vy\in V$ such that $g(\vx)=\ip{\vx,\vy}$.
\end{thm}

\begin{example}{}{}
$g(p):=\int_0^1p(x)\,\dx$ is a linear map $g:P_2(\R)\to \R$. Equip $P_2(\R)$ with the inner product for which the standard basis $\{1,x,x^2\}$ is orthonormal. Then
\[g(a+bx+cx^2)=a+\tfrac 12b+\tfrac 13c=\ip{a+bx+cx^2,1+\tfrac 12x+\tfrac 13x^2}\]
We conclude that $g(p)=\ip{p,q}$, where $q(x)=1+\frac 12x+\frac 13x^2$.
\end{example}


% By Corollary \ref{cor:orthbasis}, such a $V$ has an orthonormal basis $\beta=\{\vect{\vv}{n}\}$. Let $g:V\to \F$ be linear and define
% \[\vy:=\sum_{j=1}^n\cl{g(\vv_j)}\vv_j\]
% Certainly $\vx\mapsto \ip{\vx,\vy}$ is a linear map. Moreover,
% \[\ip{\vv_k,\vy} =\ip{\vv_k,\sum_{j=1}^n\cl{g(\vv_j)}\vv_j} =\sum_{j=1}^ng(\vv_j)\ip{\vv_k,\vv_j} =g(\vv_k)\]
% Since these linear maps are equal on a basis, they are identical. We conclude:

The idea of the proof is very simple: if $g(\vx)=\ip{\vx,\vy}$ then the nullspace of $g$ must equal $\Span\{\vy\}^\perp$\ldots

\begin{proof}
If $g$ is the zero map, take $\vy=\V0$. Otherwise, $\Rank g=1$ and
\[V=\cN(g)\oplus \cN(g)^\perp\quad\text{where}\quad\dim\cN(g)^\perp=1\hfill (\text{rank--nullity theorem and Exercise \hyperref[ex:uperpperp]{\ref*{sec:gs}.\ref*{ex:uperpperp}}})\]
Let $\vu\in \cN(g)^\perp$ be either of the two unit vectors and define, independently of $\vu$,
\[\vy:=\cl{g(\vu)}\vu\in V\]
Following the decomposition, write $\vx=\vn+\alpha\vu$ where $\vn\in\cN(g)$ and observe that
\begin{align*}
\ip{\vx,\vy}&=\ip{\vn+\alpha\vu,\cl{g(\vu)}\vu} =g(\vu)\alpha =0+g(\alpha\vu) =g(\vn+\alpha\vu) =g(\vx)
\end{align*}
The uniqueness of $\vy$ follows from the cancellation property (Lemma \ref{lemm:ipbasic}, part 4).
\end{proof}

\goodbreak



Due to the tight correspondence, the map is often decorated as $g_\vy$. Riesz's theorem indeed says that $\vy\mapsto g_\vy$ is an isomorphism $V\cong V^*$. While there are infinitely many isomorphisms between these spaces, the inner product structure identifies a \emph{canonical} or preferred choice.

\begin{cor}{}{adjexist}
Every linear map on a finite-dimensional inner product space has an adjoint.
\end{cor}



Note how only the domain is required to be finite-dimensional! Riesz's Theorem and the Corollary also apply to continuous linear operators on (infinite-dimensional) Hilbert spaces, though the proof is a little trickier.


\begin{proof}
Let $\rT\in\cL(V,W)$ where $\dim V<\infty$, and suppose $\vz\in W$ is given. Simply \emph{define} $\rT^*(\vz):=\vy$ where $\vy\in V$ is the unique vector in Riesz's Theorem arising from the linear map
\[g:V\to\F,\quad g(\vx)=\ip{\rT(\vx),\vz}\]
and check the required property:
\[\ip{\rT^*(\vz),\vx}=\ip{\vy,\vx}=\cl{g(\vx)} =\ip{\vz,\rT(\vx)}\tag*{\qedhere}\]
\end{proof}

\goodbreak


\begin{exercises}
\exstart For each inner product space $V$ and linear operator $\rT\in\cL(V)$, evaluate $\rT^*$ on the given vector.
\begin{enumerate}\setcounter{enumi}{1}
  \item[]\begin{enumerate}
    \item $V=\R^2$ with the standard inner product, $\rT\stwovec xy=\stwovec{2x+y}{x-3y}$ and $\vx=\stwovec 35$
    \item $V=\C^2$ with the standard inner product, $\rT\stwovec zw=\stwovec{2z+iw}{(1-i)z}$ and $\vx=\stwovec{3-i}{1+2i}$
    \item $V=P_1(\R)$ with $\ip{f,g}=\int_0^1 f(t)g(t)\,\dt$, $\rT(f)=f'+3f$ and $f(t)=4-2t$
  \end{enumerate}
  
	\item Suppose $A=\begin{smatrix}
	1&1\\4&3
	\end{smatrix}$ and consider the linear map $\rT=\rL_A\in\cL(\R^2)$ where $\R^2$ is equipped with the weighted inner product
	\[\ip{\vx,\vy}=4x_1y_1+x_2y_2\]
	\begin{enumerate}
	  \item Find the matrix of $\rT$ with respect to the orthonormal basis $\beta=\{\vv_1,\vv_2\}=\{\frac 12\ve_1,\ve_2\}$.
	  \item	Find the adjoint $\rT^*$ and its matrix with respect to the standard basis $\epsilon=\{\ve_1,\ve_2\}$.\smallbreak
		(\emph{Hint: the answer isn't $A^T$!})
	\end{enumerate}
	
	\item Extending Examples \ref{ex:deriv2d}, find the adjoint of $\rT=\diff x\in\cL(P_2(\R))$ with respect to:
	\begin{enumerate}
	  \item The inner product where the standard basis $\epsilon=\{1,x,x^2\}$ is orthonormal.
	  \item (Hard!)\quad The $L^2$ inner product $\int_0^1f(x)g(x)\,\dx$.
	\end{enumerate}
  
  
  \item Let $\rT(f)=f''$ be a linear transformation of $P_2(\R)$ and let $\epsilon=\{1,x,x^2\}$ be the standard basis. Find $\rT^*(a+bx+cx^2)$:
  \begin{enumerate}
    \item With respect to the inner product where $\epsilon$ is orthonormal;
    \item With respect to the $L^2$ inner product $\ip{f,g}=\int_{-1}^1f(t)g(t)\,\dt$.\smallbreak
    (\emph{Hint: $\{1,x,3x^2-1\}$ is orthogonal})
  \end{enumerate}
  
  \goodbreak

	\item Prove part 2 of Theorem \ref{thm:matrixip}.

  \item\label{ex:fundssthm} Prove the Fundamental Subspaces Theorem \ref{thm:fundssthm}.
  
  \item For each inner product space $V$ and linear transformation $g:V\to\F$, find a vector $\vy\in V$ such that $g(\vx)=\ip{\vx,\vy}$ for all $\vx\in V$.
  \begin{enumerate}
    \item $V=\R^3$ with the standard inner product, and $g\sthreevec xyz=x-2y+4z$
    \item $V=\C^2$ with the standard inner product, and $g\stwovec zw=iz-2w$
    \item $V=P_2(\R)$ with the $L^2$ inner product $\ip{f,h}=\int_0^1 f(x)h(x)\,\dx$, and $g(f)=f'(1)$
  \end{enumerate}
	
	\goodbreak
  
  
  \item\begin{enumerate}
    \item In the proof of Theorem \ref{thm:riesz}, explain why $\vy$ depends only on $g$ (not $\vu$).
    \item In the proof of Corollary \ref{cor:adjexist}, check that $g(\vx):=\ip{\rT(\vx),\vz}$ is linear.
  \end{enumerate}
  
  \item Let $\vy,\vz\in V$ be fixed vectors and define $\rT\in\cL(V)$ by $\rT(\vx)=\ip{\vx,\vy}\vz$. Show that $\rT^*$ exists and find an explicit expression.
  
  \item Suppose $A\in M_{m\times n}(\F)$. Prove that $A^*A$ is diagonal if and only if the columns of $A$ are orthogonal. What additionally would it mean if $A^*A=I$?
  
 	\item Suppose $\rT\in\cL(V)$ where $V$ is a finite-dimensional inner product space.
 	\begin{enumerate}
 	  \item Prove that the eigenvalues of $\rT^*$ are the complex conjugates of those of $\rT$.\smallbreak
 	  (\emph{Hint: relate the characteristic polynomial $p^*(t)=\det(\rT^*-t\I)$ to that of $\rT$})
 	  \item Prove that $\rT^*$ is diagonalizable if and only if $\rT$ is.
 	\end{enumerate}

	
% 	\item (\emph{Hard})\quad Let $\vu_n=(0,\ldots,0,1,0,\ldots)$ be the sequence with exactly one non-zero term in the $n\th$ position, and define $V=\Span_\R\{\vu_n\}$ as a subspace of $\ell^2$. Now define $\rT\in\cL(V)$ by
% 	\[\rT(\vx)=\left(\sum_{n=1}^\infty\frac 1nx_n\right)\vu_1\]
% 	Since $\vx\in V$, the sum only has \emph{finitely many non-zero terms} and we don't need to care about convergence; $\rT$ is therefore well-defined. If $\rT^*$ existed, show that it would have to have the formula
% 	\[\rT^*(\vy)=y_1\left(1,\frac 12,\frac 13,\ldots\right)=y_1\left(\frac 1n\right)\]
% 	Why does this \emph{not} define a map $V\to V$?\smallbreak
% 	\emph{$\rT$ is therefore a linear map without an adjoint. If we extend $\rT\in\cL(\ell^2)$ to the full space $\ell^2$, then $\rT^*$ makes perfect sense.}
	
	\item\label{exs:adjnonexist} (Hard)\quad We present two linear maps which do not have an adjoint!
	\begin{enumerate}
	  \item Since $\epsilon=\{1,x,x^2,\ldots\}$ is a basis of $P(\R)$, we may define a linear map $\rT\in\cL(P(\R))$ via $\rT(x^n)=1$ for all $n$; for instance
		\[\rT(4+3x+2x^5)=9\]
		Let $\ip{\ ,\ }$ be the inner product for which $\epsilon$ is orthonormal. If $\rT^*$ existed, show that
		\[\rT^*(1)=\sum_{n=0}^\infty x^n\]
		would be an \emph{infinite series}: $\rT^*$ therefore does not exist.
		\item For a related challenge, recall the space $\ell^2$ of square-summable real sequences. For any sequence $(x_n)_{n=1}^\infty\in\ell^2$), define $\rT\in\cL(\ell^2)$ via
		\[\rT(x_n)=\left(\sum_{n=1}^\infty\frac 1nx_n,0,0,0,0,0,\ldots\right)\]
		Find the adjoint $\rT^*$. If $V\le\ell^2$ is the subspace whose elements have only finitely many non-zero terms, show that the restriction $\at\rT V$ does not have an adjoint.
	\end{enumerate}
\end{enumerate}
\end{exercises}


\iffalse

\begin{tcolorbox}[exstyle, title={Example \ref{ex:adjointeasy}, ver.II.\quad}]
Take the same linear map $\rT=\rL_A$ as before, but this time we work with a new inner product on $\R^2$:
\[\ip{\vw,\vv}:=2xp+yq\]
Now we have a problem:
\begin{align*}
\ip{A^T\vw,\vv}
&=\ip{\begin{smatrix}
1&0\\2&3
\end{smatrix}\stwovec xy,\stwovec pq}
=2xp+(2x+3y)q\\
&\neq 2x(p+2q)+3yq
=\ip{\stwovec xy,\begin{smatrix}
1&2\\0&3
\end{smatrix}\stwovec pq}
=\ip{\vw,A\vv}
\end{align*}
The adjoint of $\rL_A$ is \emph{not} $\rL_{A^T}$. This is crucial; \emph{the adjoint of $\rT$ depends on the inner product!}\smallbreak
We can fix this via the theorem: the basis $\beta=\{\frac 1{\sqrt 2}\ve_1,\ve_2\}$ is orthonormal for the new inner product and it is easy to compute the matrix of $\rT$ with respect to $\beta$: we still have to convert back to the standard basis,
\[\left.\begin{array}{ll}
\rT(\tfrac 1{\sqrt 2}\ve_1) = \tfrac 1{\sqrt 2}\ve_1\\
\rT(\ve_2) = 2\ve_1+3\ve_2=\tfrac{2\sqrt 2}{\sqrt 2}\ve_1 +3\ve_2
\end{array}\right\}
\implies [\rT]_\beta=\begin{pmatrix}
1&2\sqrt 2\\0&3
\end{pmatrix}
\implies [\rT^*]_\beta=\begin{pmatrix}
1&0\\2\sqrt 2&3
\end{pmatrix}\]
Remember that this is with respect to $\beta$:
\[\left.\begin{array}{ll}
\rT^*(\tfrac 1{\sqrt 2}\ve_1) = \tfrac 1{\sqrt 2}\ve_1+2\sqrt 2\ve_2\\
\rT^*(\ve_2) = 3\ve_2
\end{array}\right\}
\implies [\rT^*]_\epsilon=\begin{pmatrix}
1&0\\
4&3
\end{pmatrix}
\]
As a sanity check, we repeat part of our first calculation:
\[\ip{[\rT^*]_\epsilon\vw,\vv}
=\ip{\begin{smatrix}
1&0\\4&3
\end{smatrix}\stwovec xy,\stwovec pq}
=2xp+(4x+3y)q = 2x(p+2q)+3yq
=\ip{\vw,A\vv}\]
Whew! Unfortunately the computation of adjoints is in general a gargantuan mess!
\end{tcolorbox}




\begin{example}{}{}
In Example \ref{ex:ell201} we saw that $\beta=\{f_1,f_2,f_3\}=\{1,2x-1,6x^2-6x+1\}$ is orthogonal with respect to the $L^2$ inner product $\ip{f,g}=\int_0^1f(x)g(x)\,\dx$ on $P_2(\R)$. Given $\rT=\diff x\in\cL(P_2(\R))$, we compute the adjoint in two different ways. Both rely on knowing the norms of the elements of $\beta$:
\[\Nm{f_1}^2=1,\quad \Nm{f_2}^2=\frac 13,\quad\Nm{f_3}^2=\frac 1{5}\]
\begin{enumerate}
  \item $\gamma=\{g_1,g_2,g_3\}=\{1,\sqrt{3}(2x-1),\sqrt{5}(6x^2-6x+1)\}$ is orthonormal. Observe that
  \begin{gather*}
  \rT(g_1)=0,\quad \rT(g_2)=2\sqrt 3,\qquad \rT(g_3)=6\sqrt 5(2x-1)=\frac{6\sqrt 5}{\sqrt 3}g_2=2\sqrt{15}g_2\\
  \implies [\rT]_\beta=\begin{pmatrix}
  0&2\sqrt 3&0\\
  0&0&2\sqrt{15}\\
  0&0&0
  \end{pmatrix} \implies
  [\rT^*]_\beta=\begin{pmatrix}
  0&0&0\\
  2\sqrt 3&0&0\\
  0&2\sqrt{15}&0
  \end{pmatrix}\\
  \implies \rT^*(ag_1+bg_2+cg_3)=2\sqrt 3ag_2+2\sqrt{15}bg_3
  \end{gather*}
  This can be transformed to the standard basis: you obtain the result at the bottom of the page.
  
  \item Use projections with $\beta$ directly:
  \[\rT^*(p)=\sum_{i=1}^3\frac{\ip{\rT^*(p),f_i}}{\Nm{f_i}^2}f_i =\sum_{i=1}^3\frac{\ip{p,f'_i}}{\Nm{f_i}^2}f_i =\frac{\ip{p,2}}{1/3}(2x-1) +\frac{\ip{p,12x-6}}{1/5}(6x^2-6x+1)\]
  If $p(x)$ is written with respect to $\beta$, then the inner products are easily evaluated. If you want the answer with respect to the standard basis, things are nastier. If $p(x)=a+bx+cx^2$, then
  \[\ip{p,1}=\int_0^1a+bx+cx^2\,\dx=a+\frac 12b+\frac 13c,\quad \ip{p,x}=\frac 12a+\frac 13b+\frac 14c\]
 	from which
  \begin{align*}
  \rT^*(a+bx+cx^2)&=6\left(a+\frac 12b+\frac 13c\right)(2x-1)\\
  &\qquad\qquad +30\left(a+\frac 23b+\frac 12c-a-\frac 12b-\frac 13c\right)(6x^2-6x+1)\\
  %&=(6a+3b+2c)(2x-1) +5(b+c)(6x^2-6x+1)\\
  &=(-6a+2b+3c) +(12a-24b-26c)x+(30b+30c)x^2
  \end{align*}
  With respect to the standard basis $\epsilon=\{1,x,x^2\}$, we therefore have
  \[[\rT]_\epsilon=\begin{pmatrix}
  0&1&0\\
  0&0&2\\
  0&0&0
  \end{pmatrix}\quad\text{and}\quad
  [\rT^*]_\epsilon=\begin{pmatrix}
  -6&2&3\\
  12&-24&-26\\
  0&30&30
  \end{pmatrix}\]
  This is the same result you'd obtain by applying a change of basis to the answer in version 1.
\end{enumerate}
% 
% 
% To spice things up a bit, consider an alternative inner product, where $\beta=\{1,x-1,x^2\}$ is orthonormal. With respect to this basis, we have
% \[[T(1)]_\beta=[0]_\beta=\threevec 000,\qquad [T(x-1)]_\beta=[1]_\beta=\threevec 100,\qquad [T(x^2)]_\beta=[2x]_\beta=\threevec 220\]
% It follows that the adjoint has matrix
% \[[T]_\beta^*=\begin{pmatrix}
% 0&0&0\\
% 1&0&0\\
% 2&2&0
% \end{pmatrix}\]
% In particular,
% \begin{gather*}
% T^*(1)=(x-1)+2(x^2)=2x^2+x-1,\\
% T^*(x)=T^*(x-1+1)=2x^2+(2x^2+x-1)=4x^2+x-1\\
% T^*(x^2)=0
% \end{gather*}
% Notice that this is a very different linear map from above, which is to be expected: the inner product spaces are different!
\end{example}
The moral is that adjoints are easy once you have an orthonormal basis and really hard otherwise! 

\fi

\clearpage

\subsection{Normal \& Self-Adjoint Operators and the Spectral Theorem}

We now come to the fundamental question of this chapter: for which linear operators $\rT\in\cL(V)$ can we find an orthonormal eigenbasis? Many linear maps are, of course, not even diagonalizable, so in general this is far too much to hope for! Let's see what happens if such a basis exists\ldots\medbreak

If $\beta$ is an orthonormal basis of eigenvectors of $\rT$, then
\[[\rT]_\beta=\diag(\lambda_1,\cdots,\lambda_n)\implies [\rT^*]_\beta=\diag(\cl{\lambda_1},\cdots,\cl{\lambda_n})\]
If $V$ is a \emph{real} inner product space, then these matrices are identical and so $\rT^*=\rT$. In the complex case, we instead observe that
\[[\rT\rT^*]_\beta=\diag(\nm{\lambda_1}^2,\cdots,\nm{\lambda_n}^2)=[\rT^*\rT]_\beta \implies \rT\rT^*=\rT^*\rT\]

\begin{defn}{}{selfadj}
Suppose $\rT$ is a linear operator on an inner product space $V$ and assume $\rT$ has an adjoint. We say that $\rT$ is,
\begin{itemize}
	\item[]\emph{Normal} if $\rT\rT^*=\rT^*\rT$,
	\item[]\emph{Self-adjoint} if $\rT^*=\rT$.
\end{itemize}
The definitions for square matrices over $\R$ and $\C$ are identical, where $^*$ now denotes the conjugate-transpose.
\end{defn}

A real self-adjoint matrix $A\in M_n(\R)$ is plainly \emph{symmetric}: $A^T=A$. A complex self-adjoint matrix is also called \emph{Hermitian}: $A^*=A$.\smallbreak
If $\rT$ is self-adjoint then it is certainly normal, but the converse is false as the next example shows.

\begin{example}{}{rotatenotselfadj}
The (non-symmetric) real matrix $A=\begin{smatrix}
2&-1\\
1&2
\end{smatrix}$ is normal but not self-adjoint:
\[AA^T=\begin{pmatrix}
2&-1\\
1&2
\end{pmatrix}\begin{pmatrix}
2&1\\
-1&2
\end{pmatrix}=\begin{pmatrix}
5&0\\
0&5
\end{pmatrix}=\begin{pmatrix}
2&1\\
-1&2
\end{pmatrix}\begin{pmatrix}
2&-1\\
1&2
\end{pmatrix}=A^TA\]
\end{example}

More generally, every non-zero \emph{skew-hermitian} matrix $A^*=-A$ is normal but \emph{not} self-adjoint:
\[A^*=-A\implies AA^*=-A^2=A^*A\]

We saw above that linear maps with an orthonormal eigenbasis are either self-adjoint or normal depending whether the inner product space is real or complex. Amazingly, this provides a complete characterisation of such maps!


\begin{thm}{Spectral Theorem, version 1}{spec}
Let $\rT$ be a linear operator on a finite-dimensional inner product space $V$.
\begin{description}
  \item[\normalfont\emph{Complex case}:] $V$ has an orthonormal basis of eigenvectors of $\rT$ if and only if $\rT$ is normal.
  \item[\normalfont\emph{Real case}:] $V$ has an orthonormal basis of eigenvectors of $\rT$ if and only if $\rT$ is self-adjoint.
\end{description}
\end{thm}

The theorem gets is name from the \emph{spectrum} (set of eigenvalues) of $\rT$.


\goodbreak

\begin{examples}{}{specthm}
\exstart We diagonalize the self-adjoint linear map $\rT=\rL_A\in\cL(\R^2)$ where $A=\begin{smatrix}
	6&3\\3&-2
	\end{smatrix}$.
	\[\def\arraystretch{1.2}\begin{array}{@{}ll}
		\text{\emph{Characteristic polynomial}} &p(t)=(6-t)(-2-t)-9=t^2-4t-21=(t-7)(t+3)\\
		\text{\emph{Eigenvalues}} &\lambda_1=7,\quad \lambda_2=-3\\
		\text{\emph{Eigenvectors (normalized)}} &\vw_1=\frac 1{\sqrt{10}}\stwovec{3}{1},\quad \vw_2=\frac 1{\sqrt{10}}\stwovec{-1}{3}
	\end{array}\]
\begin{enumerate}\setcounter{enumi}{1}
	\item[]	The basis $\beta=\{\vw_1,\vw_2\}$ is orthonormal, with respect to which $[\rT]_\beta=\begin{smatrix}
	7&0\\0&-3
	\end{smatrix}$ is diagonal.
	
  \item\label{ex:schur1} The map $\rT=\rL_A\in\cL(\R^2)$ where $A=\begin{smatrix}
	1&3\\0&-2
	\end{smatrix}$ is neither self-adjoint nor normal:
	\[AA^*=\begin{pmatrix}
	1&3\\0&-2
	\end{pmatrix}\begin{pmatrix}
	1&0\\3&-2
	\end{pmatrix}
	=
	\begin{pmatrix}
	10&-6\\-6&4
	\end{pmatrix}
	\neq \begin{pmatrix}
	1&3\\3&13
	\end{pmatrix}
	=
	\begin{pmatrix}
	1&0\\3&-2
	\end{pmatrix}\begin{pmatrix}
	1&3\\0&-2
	\end{pmatrix}
	=A^*A\]
	It is diagonalizable, indeed
	\[\gamma=\left\{\twovec 10,\twovec 3{-1}\right\} \rightsquigarrow [\rT]_\gamma=\begin{pmatrix}
	1&0\\0&-2
	\end{pmatrix}\]
	In accordance with the spectral theorem, $\gamma$ is not orthogonal.
	
	\item Let $A=\begin{smatrix}0&-1\\1&0\end{smatrix}$ and consider $\rT=\rL_A$ acting on both $\C^2$ and $\R^2$. Since $\rT$ is normal but not self-adjoint, we'll see how the field really matters in the spectral theorem.\smallbreak
	First the complex case: $\rT\in\cL(\C^2)$ is normal and thus diagonalizable with respect to an orthonormal basis of eigenvectors. Here are the details.
	\[\def\arraystretch{1.2}\begin{array}{@{}ll}
		\text{\emph{Characteristic polynomial}} &p(t)=t^2+1=(t-i)(t+i)\\
		\text{\emph{Eigenvalues}} &\lambda_1=i,\quad \lambda_2=-i\\
		\text{\emph{Eigenvectors (normalized)}} &\vw_1=\frac 1{\sqrt{2}}\stwovec{i}{1},\quad \vw_2=\frac 1{\sqrt{2}}\stwovec{-i}{1}
	\end{array}\]
	Certainly $\ip{\vw_1,\vw_2}=0$, $\beta=\{\vw_1,\vw_2\}$ is orthonormal, and $[\rT]_\beta=\begin{smatrix}
	i&0\\0&-i
	\end{smatrix}$ is diagonal.\smallbreak	
	Now for the real case: $\rT\in\cL(\R^2)$ is not self-adjoint and thus should not be diagonalizable with respect to an orthonormal basis of eigenvectors. Indeed this is trivial; the characteristic polynomial has no roots in $\R$ and so there are no real eigenvalues! It is also clear geometrically: $\rT$ is rotation by \ang{90} counter-clockwise around the origin, so it has no eigenvectors.
	
	\item Let $A=\begin{smatrix}3&i\\-i&3\end{smatrix}$ and consider the self-adjoint operator $\rT=\rL_A\in\cL(\C^2)$.
	\[\def\arraystretch{1.2}\begin{array}{@{}ll}
		\text{\emph{Characteristic polynomial}} &p(t)=t^2-6t+9-1=(t-2)(t-4)\\
		\text{\emph{Eigenvalues}} &\lambda_1=2,\quad \lambda_2=4\\
		\text{\emph{Eigenvectors (normalized)}} &\vw_1=\frac 1{\sqrt{2}}\stwovec{1}{i},\quad \vw_2=\frac 1{\sqrt{2}}\stwovec{1}{-i}
	\end{array}\]
	With respect to the orthonormal basis $\beta=\{\vw_1,\vw_2\}$, we have $[\rT]_\beta=\begin{smatrix}
	2&0\\0&4
	\end{smatrix}$.
\end{enumerate}
\end{examples}
\goodbreak

\boldsubsubsection{Proving the Spectral Theorem for Self-Adjoint Operators}


\begin{lemm}{Basic properties of self-adjoint operators}{selfadjbasic}
Let $\rT\in\cL(V)$ be self-adjoint.
\begin{enumerate}\itemsep2pt
  \item If $W\le V$ is $\rT$-invariant then the restriction $\at\rT W$ is self-adjoint;
  \item Every eigenvalue of $\rT$ is real;
  \item If $\dim V$ is finite then $\rT$ has an eigenvalue.
\end{enumerate}
\end{lemm}

It is irrelevant whether $V$ is real or complex. The previous example demonstrates part 2; even when $V=\C^2$ is a complex inner product space, the eigenvalues of a self-adjoint matrix are real.

\begin{proof}
\begin{enumerate}\itemsep2pt
  \item Let $\vw_1,\vw_2\in W$. Then
  \[\ip{\at\rT W(\vw_1),\vw_2}= \ip{\rT(\vw_1),\vw_2}= \ip{\vw_1,\rT(\vw_2)} =\ip{\vw_1,\at\rT W(\vw_2)}\]
  \item Suppose $(\lambda,\vx)$ is an eigenpair. Then
  \[\lambda\Nm{\vx}^2=\ip{\rT(\vx),\vx}=\ip{\vx,\rT(\vx)}=\cl\lambda\Nm{\vx}^2\implies \lambda\in\R\]
  \item This is trivial if $V$ is complex since every characteristic polynomial splits over $\C$. We therefore assume $V$ is real. Choose any orthonormal basis $\gamma$ of $V$, let $A=[\rT]_\gamma\in M_n(\R)$, and define $\rS:=\rL_A\in\cL(\C^{n})$. Then;
	\begin{itemize}\itemsep2pt
	  \item The characteristic polynomial of $\rS$ splits over $\C$, whence there exists an eigenvalue $\lambda\in\C$.
  	\item The characteristic polynomials of $\rS$ and $\rT$ are \emph{identical} (to that of $A$).
  	\item $\rS$ is self-adjoint and thus (part 2) $\lambda\in\R$.
	\end{itemize}
  It follows that $\rT$ has the same \emph{real} eigenvalue $\lambda$.\qedhere
\end{enumerate}
\end{proof}

We are now able to prove the spectral theorem for self-adjoint operators on a finite-dimensional inner product space $V$. The argument applies regardless of whether $V$ is real or complex.

\begin{proof}[Proof of the Spectral Theorem (self-adjoint case)]
We prove by induction on $\dim V$.\smallbreak
(Base case)\quad If $\dim V=1$, then $V=\Span\{\vx\}$ and $\rT(\vx)=\lambda\vx$ for some unit vector $\vx$ and scalar $\lambda\in\R$; plainly $\{\vx\}$ is an orthonormal eigenbasis for $\rT$.\smallbreak
(Induction step)\quad Fix $n\in\N$ and assume that \emph{every} self-adjoint operator on \emph{every} inner product space of dimension $n$ satisfies the spectral theorem. Let $\dim V=n+1$ and $\rT\in\cL(V)$ be self-adjoint.\smallbreak
By part 3 of the Lemma, $\rT$ has an eigenpair $(\lambda,\vx)$ where we may assume $\vx$ has unit length. Let $W=\Span\{\vx\}^\perp$. If $\vw\in W$, then
\[\ip{\vx,\rT(\vw)}=\ip{\rT(\vx),\vw}=\lambda\ip{\vx,\vw}=0\tag{$\ast$}\]
whence $W$ is $\rT$-invariant. Plainly $\dim W=n$. By part 1 of the Lemma, $\at\rT W$ is self-adjoint. By the induction hypothesis, $\at\rT{W}$ is diagonalized by some orthonormal basis $\gamma$ of $W$. But then  $\rT$ is diagonalized by the orthonormal basis $\beta=\gamma\cup\{\vx\}$ of $V$.
\end{proof}

\goodbreak

\subsubsection*{Proving the Spectral Theorem for Normal Operators}

What changes for normal operators on complex inner product spaces? Not much! Indeed the proof is almost identical when $\rT$ is merely normal.
\begin{itemize}
  \item We don't need parts 2 and 3 of Lemma \ref{lemm:selfadjbasic}: every linear operator on a finite-dimensional complex inner product space has an eigenvalue and we no longer care whether eigenvalues are real.
  \item Two parts of the induction step need fixed:
  \begin{itemize}
    \item $W$ being $\rT$-invariant: This isn't quite as simple as ($\ast$), but thankfully part 3 of the next result provides the needed correction.
    \item $\at\rT W$ being normal: We need a replacement for part 1 of Lemma \ref{lemm:selfadjbasic}; this is a little more involved.
	\end{itemize}
\end{itemize}

Rather than write out all the details, we leave this to Exercises \ref{ex:normalspec1} and \ref{ex:normalspec2}.\smallbreak

For completeness, and as an analogue/extension of Lemma \ref{lemm:selfadjbasic}, we summarize some of the basic properties of normal operators. These also apply to self-adjoint operators as a special case.

\begin{lemm}{Basic properties of normal operators}{normalbasic}
Let $\rT$ be normal on $V$. Then:
\begin{enumerate}
  \item $\forall \vx\in V$, $\Nm{\rT(\vx)}=\Nm{\rT^*(\vx)}$.
  \item $\rT-t\I$ is normal for any scalar $t$.
  \item $\rT(\vx)=\lambda\vx\iff \rT^*(\vx)=\cl\lambda\vx$ so that $\rT$ and $\rT^*$ have the same eigenvectors and conjugate eigenvalues. This recovers the previously established fact that $\lambda\in\R$ if $\rT$ is self-adjoint.
  \item Distinct eigenvalues of $\rT$ have orthogonal eigenvectors.
\end{enumerate}
\end{lemm}


\begin{proof}
\begin{enumerate}
  \item $\Nm{\rT(\vx)}^2=\ip{\rT(\vx),\rT(\vx)}=\ip{\rT^*\rT(\vx),\vx} =\ip{\rT\rT^*(\vx),\vx}=\ip{\rT^*(\vx),\rT^*(\vx)}=\Nm{\rT^*(\vx)}^2$.
  \item $\ip{\vx,(\rT-t\I)(\vy)}=\ip{\vx,\rT(\vy)}-\cl t\ip{\vx,\vy} =\ip{\rT^*(\vx),\vy}-\ip{\cl t\vx,\vy} =\ip{(\rT^*-\cl t\I)(\vx),\vy}$ shows that $\rT-t\I$ has adjoint $\rT^*-\cl t\I$. It is trivial to check that these commute.
  \item $\rT(\vx)=\lambda\vx\iff \Nm{(\rT-\lambda \I)(\vx)}=0\underset{1 \& 2}{\overset{\text{parts}}{\iff}} \Nm{(\rT^*-\cl\lambda \I)(\vx)}=0 \iff \rT^*(\vx)=\cl\lambda\vx$.
  \item In part this follows from the spectral theorem, but we can also prove more straightforwardly. Suppose $\rT(\vx)=\lambda\vx$ and $\rT(\vy)=\mu\vy$ where $\lambda\neq \mu$. By part 3,
  \begin{align*}
  \lambda\ip{\vx,\vy} &=\ip{\lambda\vx,\vy} =\ip{\rT(\vx),\vy} =\ip{\vx,\rT^*(\vy)} = \ip{\vx,\cl\mu\vy}\\
    & =\mu\ip{\vx,\vy}
  \end{align*}
  This is a contradiction unless $\ip{\vx,\vy}=0$.\hfill\qedhere
\end{enumerate}
\end{proof}
\vfil\vfil
\goodbreak

\boldsubsubsection{Schur's Lemma}

It is reasonable to ask how useful an orthonormal basis can be in general. Here is one answer.

% \begin{example}[breakable]{}{}
% Consider $T=L_A$ where $A=\begin{smatrix}
% 0&3+4i\\-3+4i&0
% \end{smatrix}$. This is skew-hermitian ($A^*=-A$) and thus normal but not self-adjoint. To find its eigenpairs, we solve
% \[p(t)=t^2-(3+4i)(-3+4i)=t^2+25=0\implies t=\pm 5i\]
% The eigenvectors can then be found:
% \[\lambda=5i\implies \begin{pmatrix}
%   -5i&3+4i\\
%   -3+4i&-5i
%   \end{pmatrix}\vw=\V0\implies \text{choose }\vw_1=\twovec{3+4i}{5i}\]
%   while the other eigenpair is seen to be $\cl\lambda,\vw_2=\stwovec{3+4i}{-5i}$. Observe that these eigenvectors are orthogonal:
% \[\ip{\vw_1,\vw_2}=\vw_2^*\vw_1=(\cl{3+4i})(3+4i)+(\cl{-5i})(5i)= 25-25=0\]
% Note also that the adjoint $A^*=-A$ has the same eigenvectors but with conjugate eigenvalues:
% \[A^*\vw=-A\vw=-\lambda\vw=\cl\lambda\vw\]
% \end{example}
% 
% As a consequence of Lemma \ref{lemm:normalbasic}, any \emph{diagonalizable} normal or self-adjoint operators necessarily have an orthogonal eigenbasis. The hard part of the spectral theorem is proving that such operators are diagonalizable in the first place! This can be done in one shot, but first we develop a useful and interesting result.

\begin{lemm}{Schur}{schur}
Suppose $\rT$ is a linear operator on a finite-dimensional inner product space $V$. If the characteristic polynomial of $\rT$ splits, then there exists an orthonormal basis $\beta$ of $V$ such that $[\rT]_\beta$ is upper-triangular.
\end{lemm}

The spectral theorem is a special case; since the proof is similar, we leave it to the exercises.\smallbreak

The conclusion of Schur's lemma is weaker than the spectral theorem, though it applies to more operators: indeed if $V$ is complex, it applies to \emph{any} $\rT$! Every example of the spectral theorem is also an example of Schur's lemma. Example \hyperref[ex:schur1]{\ref*{ex:specthm}.\ref*{ex:schur1}} provides another, since the matrix $A$ is already upper triangular with respect to the standard orthonormal basis. Here is a another example.

\begin{example}{}{}
Consider $\rT(f)=2f'(x)+xf(1)$ as a linear map $\rT\in\cL(P_1(\R))$ with respect to the $L^2$ inner product $\ip{f,g}=\int_0^1f(t)g(t)\,\dt$. We have
\[\rT(a+bx)=2b+(a+b)x\]
If $[\rT]_\beta$ is to be upper-triangular, the first vector in $\beta$ must be an eigenvector of $\rT$.
It is easily checked that $f_1=1+x$ is such with eigenvalue 2. To find a basis satisfying Schur's lemma, we need only find $f_2$ orthogonal to this and then normalize. This can be done by brute force since the problem is small, but for the sake of practice we apply Gram--Schmidt to the polynomial 1:
\[1-\frac{\ip{1,1+x}}{\Nm{1+x}^2}(1+x)= 1-\frac{1+\frac 12}{1+1+\frac 13}(1+x)=\frac 1{14}(5-9x) \implies f_2=5-9x\]
Indeed we obtain an upper-triangular matrix for $\rT$:
\[\rT(f_2)=-18-4x =-13(1+x)-(5-9x) =-13f_1-f_2\implies [\rT]_{\{f_1,f_2\}}=\begin{pmatrix}
2&-13\\
0&-1
\end{pmatrix}\]
We can also work with the corresponding orthonormal basis as posited in the theorem, though the matrix is messier:
\[\beta=\{g_1,g_2\}=\left\{\sqrt{\frac 37}(1+x),\frac 1{\sqrt 7}(5-9x)\right\} \implies [\rT]_\beta=\begin{pmatrix}
2&-\frac{13}{\sqrt 3}\\
0&-1
\end{pmatrix}\]
Alternatively, we could have started with the other eigenvector $h_1=2-x$: an orthogonal vector to this is $h_2=4-9x$, with respect to which
\[[\rT]_{\{h_1,h_2\}}=\begin{pmatrix}
-1&-13\\
0&2
\end{pmatrix}\]
In both cases the eigenvalues are down the diagonal, as must be for an upper-triangular matrix.
\end{example}

In general, it is difficult to quickly find a suitable basis satisfying Schur's lemma. After trying the proof in the exercises, you should be able to describe a method, though it is impractically slow!

\goodbreak

% \parbox{\linewidth}{Unfortunately the tactic of extending a single eigenvector to a basis doesn't work so well when $\dim V\ge 3$. Instead we appeal to the more general idea of an invariant subspace. Recall that $U\le V$ being $T$-invariant means that $T(U)\le U$: moreover, if $\gamma$ is a basis of $U$ extended to a basis $\beta=\gamma\cup\eta$ of $V$, then the matrix of $T$ with respect to $\beta$ has block form
% \[[T]_\beta=\left(\begin{array}{c|c}
% A&B\\\hline
% 0&C
% \end{array}\right)\]
% The core of the proof is an induction argument: if $\dim U=\dim V-1$ and $A$ is upper-triangular, then $[T]_\beta$ is upper-triangular.\footnote{In the proof, $U$ is the subspace $W^\perp$, and $\eta=\{\vu\}$.}}


% \begin{proof}[Proof of Schur's Lemma]
% If $\dim V=1$, then the result is trivial.\\[5pt]
% Now suppose the result is true for all inner product spaces with dimension $n$. Suppose also that $\dim V=n+1$, and that the characteristic polynomial of $\rT\in \cL(V)$ splits. It follows that $\rT$ has at least one eigenvalue $\mu$ and that its adjoint has an eigenpair $(\cl\mu,\vu)$. Let $W=\Span\{\vu\}^\perp$ and let $\vw\in W$. Then
% \[\ip{\rT(\vw),\vu}=\ip{\vw,\rT^*(\vu)}=\mu\ip{\vw,\vu}=0\implies \rT(W)\le W\]
% whence $W$ is $\rT$-invariant. Recall chapter 5: the characteristic polynomial of the restriction $\at\rT W$ divides that of $\rT$ and therefore splits. By the induction hypothesis, $W$ has a basis $\gamma$ with respect to which $[\at\rT{W}]_\gamma$ is upper-triangular. Let $\beta=\gamma\cup\{\vu\}$ to finish the induction and thus the proof.
% \end{proof}

% We are now (finally!) able to prove the spectral theorem. First we tackle the real case. Since the proof that follows is a little sneaky, we first give a synopsis. To invoke Schur's Lemma we need to see that the characteristic polynomial of $T$ (an operator on a \emph{real} inner product space) splits over $\R$. The trick is to first transfer to a complex inner product space so we can appeal to the Fundamental Theorem of Algebra. Here are the main steps:
% \begin{enumerate}
%   \item Create an associated self-adjoint linear operator $S$ on $\C^n$ with the \emph{same} characteristic polynomial as $T$.
%   \item The Fundamental Theorem of Algebra says this splits over $\C$.
%   \item Self-adjointness says the eigenvalues are \emph{real}.
%   \item Apply Schur's Lemma for an orthonormal basis.
%   \item Show that the resulting matrix is diagonal.
% \end{enumerate}
% 
% \begin{proof}[Proof of the Spectral Theorem for real inner product spaces]
% Let $\dim V=n$, $T$ be self-adjoint, choose any orthonormal basis $\gamma$ of $V$, let $A=[T]_\gamma$ and define $S=L_A\in\cL(\C^n)$ to be multiplication by $A$. With respect to the standard basis $\epsilon$ of $\C^n$, the matrix of $S$ is $[S]_\epsilon=A$. Two facts are immediate:
% \begin{itemize}
%   \item The characteristic polynomials of $S$ and $T$ are identical (to that of $A$).
%   \item $S$ is self-adjoint and thus (Lemma \ref{lemm:normalbasic}) has real eigenvalues.
% \end{itemize}
% By the Fundamental Theorem of Algebra, the characteristic polynomial splits \emph{over $\C$}
% \[p(t)=(\lambda_1-t)\cdots(\lambda_n-t)\]
% where the $\lambda_j$ are the \emph{real} eigenvalues of $S$. It follows that $p(t)$ splits over $\R$.\\[5pt]
% By Schur's Lemma, there is an orthonormal basis $\beta$ of $V$ with respect to which $B:=[T]_\beta$ is upper-triangular. However (Lemma \ref{lemm:adjmatrix}), since $T$ is self-adjoint and $\beta$ is orthonormal,
% \[B^*=[T^*]_\beta=[T]_\beta=B\]
% Since $B$ is simultaneously upper- and lower-triangular, it must be diagonal.\\[5pt]
% For the converse, suppose $\beta$ is an orthonormal basis of eigenvectors of $T$. Then $[T]_\beta$ is diagonal with \emph{real} entries. It follows that $[T^*]_\beta=[T]_\beta^*=[T]_\beta$, whence $T^*=T$ is self-adjoint.
% \end{proof}
% 
% The proof for the complex case is quite different. We can instantly apply the Fundamental Theorem of Algebra to invoke Schur and obtain a suitable basis. Since normality is a weaker condition than self-adjointness, the difficulty lies in showing that this basis consists entirely of eigenvectors.
% 
% \begin{proof}[Proof of Spectral Theorem for complex inner product spaces]
% Let $\dim V=n$ and $T$ be normal. By Schur's Lemma, there exists an orthonormal basis $\beta=\{\vect{\vw}{n}\}$ for which $A:=[T]_\beta$ is upper-triangular. We prove by induction that $\beta$ is a basis of eigenvectors.\\[5pt]
% For the base case, $\vw_1$ is an eigenvector by the upper-triangularity of $A$.\\[5pt]
% Now suppose $\vect{\vw}{k-1}$ are eigenvectors, with $T(\vw_j)=\lambda_j\vw_j$. By Lemma \ref{lemm:normalbasic} $T^*(\vw_j)=\cl{\lambda_j}\vw_j$. By Corollary \ref{cor:matrixip} and the orthonormality of $\beta$,
% \[j<k\implies A_{jk}=\ip{T(\vw_k),\vw_j}=\ip{\vw_k,T^*(\vw_j)} =\lambda_j\ip{\vw_k,\vw_j}=0\]
% However, since $A$ is upper-triangular,
% \[T(\vw_{k})=\sum_{j=1}^{k}A_{jk}\vw_j =A_{kk}\vw_k\]
% whence $\vw_k$ is an eigenvector of $T$.
% By induction, $\beta$ is a basis of eigenvectors and $A$ is diagonal.\\[5pt]
% Conversely, if $\beta$ is an orthonormal basis of eigenvectors, then $[T]_\beta$ is diagonal and so $[T^*]_\beta=[T]_\beta^*$ is diagonal. Since diagonal matrices commute, we see that $T^*T=TT^*$: that is, $T$ is normal.
% \end{proof}



 

\begin{exercises}
\exstart For each linear operator $\rT$ on an inner product space $V$, decide whether $\rT$ is normal, self-adjoint, or neither. If the spectral theorem permits, find an orthonormal eigenbasis.
\begin{enumerate}\setcounter{enumi}{1}
  \item[]\begin{enumerate}
    \item $V=\R^2$ and $\rT\stwovec xy=\stwovec{2a-b}{-2a+5b}$
    \item $V=\R^3$ and $\rT\sthreevec xyz=\sthreevec{-a-b}{5b}{4a-2b+5c}$
    \item $V=\C^2$ and $\rT\stwovec zw=\stwovec{2z+iw}{z+2w}$
    \item $V=\R^4$ with $\rT:(\ve_1,\ve_2,\ve_3,\ve_4)\mapsto (\ve_3,\ve_4,\ve_1,\ve_2)$
    \item $V=P_2(\R)$ with $\ip{f,g}=\int_0^1f(t)g(t)\,\dt$ and $\rT(f)=f'$\\[2pt]
    (\emph{Hint: Don't compute $\rT^*$! Instead assume $\rT$ is normal and aim for a contradiction\ldots})
  \end{enumerate}
  
  \item Let $\rT(f(x))=f'(x)+4xf(0)$ where $\rT\in\cL(P_1(\R))$ and $\ip{f,g}=\int_{-1}^1f(t)g(t)\,\dt$. Find an orthonormal basis of $P_1(\R)$ with respect to which the matrix of $\rT$ is upper-triangular.
  
  \item Suppose $\rS,\rT$ are self-adjoint operators on an inner product space $V$. Prove that $\rS\rT$ is self-adjoint if and only if $\rS\rT=\rT\rS$.\\[2pt]
  (\emph{Hint: recall Theorem \ref{thm:matrixip}})
  
  
  %\item Let $T$ be a normal operator on a finite-dimensional inner product space $V$ and let $W\le V$. Prove that if $W$ is $T$-invariant then $W$ is also $T^*$-invariant.
  
  \item Let $\rT$ be normal on a finite-dimensional inner product space $V$. Prove that $\cN(\rT^*)=\cN(\rT)$ and that $\cR(\rT^*)=\cR(\rT)$.\\[2pt]
  (\emph{Hint: Use Lemma \ref{lemm:normalbasic} and the Fundamental Subspaces Theorem \ref{thm:fundssthm}})
  
  \item Let $\rT$ be self-adjoint on a finite-dimensional inner product space $V$. Prove that
  \[\forall\vx\in V,\quad \Nm{\rT(\vx)\pm i\vx}^2=\Nm{\rT(\vx)}^2+\Nm{\vx}^2\]
  Hence prove that $\rT-i\I$ is invertible and that $[(\rT-i\I)^{-1}]^*=(\rT+i\I)^{-1}$.
  
    
  \item\label{ex:normalspec1} Let $W$ be a $\rT$-invariant subspace of an inner product space $V$ and let $\at\rT W\in\cL(W)$ be the restriction of $\rT$ to $W$. Prove:
  \begin{enumerate}
    %\item $T$ self-adjoint $\implies \at TW$ is self-adjoint.
    \item $W^\perp$ is $\rT^*$-invariant.
    \item If $W$ is both $\rT$- and $\rT^*$-invariant, then $(\at\rT W)^*=\at{(\rT^*)}W$.
    \item If $W$ is both $\rT$- and $\rT^*$-invariant and $\rT$ is normal, then $\at\rT W$ is normal.
  \end{enumerate}
  
  \item\label{ex:normalspec2} Use the previous question to complete the proof of the spectral theorem for a normal operator on a finite-dimensional complex inner product space.
  
  \item\begin{enumerate}
    \item Suppose $\rS$ is a normal operator on a finite-dimensional \emph{complex} inner product space, all of whose eigenvalues are \emph{real.} Prove that $\rS$ is self-adjoint.
    \item Let $\rT$ be a normal operator on a finite-dimensional \emph{real} inner product space $V$ whose characteristic polynomial splits. Prove that $\rT$ is self-adjoint and that there exists an orthonormal basis of $V$ of eigenvectors of $\rT$.\\[2pt]
    (\emph{Hint: Mimic the proof of Lemma \ref{lemm:selfadjbasic} part 3 and use part (a)})
  \end{enumerate}
  
  \item Prove Schur's lemma by induction, similarly to the proof of the spectral theorem.\\[2pt]
  (\emph{Hint: $\rT^*$ has an eigenvector $\vx$; why? Now show that $W=\Span\{\vx\}^\perp$ is $\rT$-invariant\ldots})
\end{enumerate}
\end{exercises}

\clearpage

\subsection{Unitary and Orthogonal Operators and their Matrices}

% In this section considering the diagonalization of a normal or self-adjoint matrix in accordance with the spectral theorem. Here is an example:
% 
% \begin{example}{}{}
% Consider the self-adjoint (symmetric) matrix $A=\begin{smatrix}
% 7&12\\12&0
% \end{smatrix}=A^T$. This is easily seen to have orthogonal eigenbasis $\beta=\left\{\frac 15\stwovec 43,\frac 15\stwovec{-3}4\right\}$ and may be diagonalized as
% \[A=\frac 15\begin{pmatrix}
% 4&-3\\
% 3&4
% \end{pmatrix}
% \begin{pmatrix}
% 16&0\\
% 0&-9
% \end{pmatrix}
% \left[\frac 15\begin{pmatrix}
% 4&-3\\
% 3&4
% \end{pmatrix}\right]^{-1}
% =\frac 15\begin{pmatrix}
% 4&-3\\
% 3&4
% \end{pmatrix}
% \begin{pmatrix}
% 16&0\\
% 0&-9
% \end{pmatrix}
% \left[\frac 15\begin{pmatrix}
% 4&3\\
% -3&4
% \end{pmatrix}\right]\]
% Note the decomposition: writing $A=UDU^{-1}$, we see that $U^{-1}=U^T$.
% \end{example}
% The same thing happens in general.
% 
% A normal matrix $A\in M_n(\C)$, has an orthonormal eigenbasis $\beta=\{\vect{\vu}{n}\}$: thus
% \[A=UDU^{-1}\quad\text{where}\quad U=(\vu_1\ \cdots\ \vu_n)\quad\text{and}\quad D=\diag(\lambda_1,\ldots,\lambda_n)\] 
% Here $U$ is the matrix whose $j^\text{th}$ column is $\vu_j$ and $D$ is the matrix with the eigenvalues of $A$ down the diagonal. Considering the conjugate-transpose $\vu_j^*$ as a row vector, we see that
% \[U^*U=\threevec{\vu_1^*}{\vdots}{\vu_n^*}(\vu_1\ \cdots\ \vu_n) =\begin{pmatrix}
% \vu_1^*\vu_1&\cdots&\vu_1^*\vu_n\\
% \vdots&&\vdots\\
% \vu_n^*\vu_1&\cdots&\vu_n^*\vu_n
% \end{pmatrix} =I\quad\text{since}\quad \vu_k^*\vu_j=\ip{\vu_j,\vu_k}\]

In this section we focus on length-preserving transformations of an inner product space.

\begin{defn}{}{}
A \emph{linear\footnotemark isometry} of an inner product space $V$ is a linear map $\rT$ satisfying
\[\forall \vx\in V,\ \Nm{\rT(\vx)}=\Nm{\vx}\]
\end{defn}

\footnotetext{There also exist \emph{non-linear} isometries: for instance \emph{translations} ($\rT(\vx)=\vx+\va$ for any constant $\va$) and \emph{complex conjugation} ($\rT(\vx)=\cl\vx$) on $\C^n$. Together with linear isometries, these essentially comprise all isometries in finite dimensions.}

Every eigenvalue of an isometry must have modulus 1: if $\rT(\vw)=\lambda\vw$, then
\[\Nm{\vw}^2=\Nm{\rT(\vw)}^2=\Nm{\lambda\vw}^2=\nm\lambda^2\Nm{\vw}^2\]

\begin{example}{}{easyrotate}
Let $\rT=\rL_A\in\cL(\R^2)$, where $A=\frac 15\begin{smatrix}
4&-3\\3&4
\end{smatrix}$. Then
\[\Nm{\rT\twovec xy}^2=\Nm{\frac 15\twovec{4x-3y}{3x+4y}}^2=\frac 1{25}\left((4x-3y)^2+(3x+4y)^2\right)=x^2+y^2=\Nm{\twovec xy}^2\]
This matrix is very special in that its inverse equals its transpose:
\[A^{-1}=\frac 1{\frac{16}{25}+\frac 9{25}}\begin{pmatrix}
4&3\\-3&4
\end{pmatrix} =\frac 1{5}\begin{pmatrix}
4&3\\-3&4
\end{pmatrix}=A^T\]
\end{example}

We call such matrices \emph{orthogonal.} The simple version of what follows is that every linear isometry on $\R^n$ is multiplication by an orthogonal matrix.

\begin{defn}{}{unitary}
A \emph{unitary operator} $\rT$ on an inner product space $V$ is an invertible linear map satisfying $\rT^*\rT=\I=\rT\rT^*$. A \emph{unitary matrix} is a (real or complex) matrix satisfying $A^*A=I$.\smallbreak
If $V$ is real, we usually call these \emph{orthogonal operators/matrices}; this isn't necessary, since \emph{unitary} encompasses both real and complex spaces. An orthogonal matrix satisfies $A^TA=I$.
\end{defn}

\begin{example}{}{}
The matrix $A=\frac 13\begin{smatrix}
i&2+2i\\2-2i&i
\end{smatrix}$ is unitary:
\[A^*A=\frac 19\begin{pmatrix}
-i&2+2i\\2-2i&-i
\end{pmatrix}\begin{pmatrix}
i&2+2i\\2-2i&i
\end{pmatrix}
=\frac 19\begin{pmatrix}
-i^2+4+4&(-i+i)(2+2i)\\(i-i)(2-2i)&4+4-i^2
\end{pmatrix}
=I\]
\end{example}

If $V$ is finite-dimensional, the operator/matrix notions correspond straightforwardly. By Theorem \ref{thm:matrixip}, if we choose any orthonormal basis $\beta$ of $V$, then
\[\rT\in\cL(V)\text{ is unitary/orthogonal } \iff [\rT]_\beta \text{ is unitary/orthogonal}\]
Moreover, we need only assume $\rT^*\rT=\I$ (or $\rT\rT^*=\I$) if $V$ is finite-dimensional: if $\beta$ is an orthonormal basis, then
\[\rT^*\rT=\I\iff[\rT^*]_\beta[\rT]_\beta=I\iff [\rT]_\beta[\rT^*]_\beta=I\iff \rT\rT^*=\I\]
In infinite dimensions, we need $\rT^*$ to be both the left- and right-inverse of $\rT$. This isn't an empty requirement (see Exercise \ref{ex:isometrynotunitary}).


\goodbreak

We now tackle the correspondence between unitary operators and isometries.

\begin{thm}{}{}
Let $\rT$ be a linear operator on an inner product space $V$.
\begin{enumerate}
  \item If $\rT$ is a unitary/orthogonal operator, then it is a linear isometry.
  \item If $\rT$ is a linear isometry and $V$ is finite-dimensional, then $\rT$ is unitary/orthogonal.
\end{enumerate}
\end{thm}

\begin{proof}
\begin{enumerate}
  \item If $\rT$ is unitary, then
	\[\forall\vx,\vy\in V,\quad\ip{\vx,\vy}=\ip{\rT^*\rT(\vx),\vy}=\ip{\rT(\vx),\rT(\vy)}\tag{$\dag$}\]
	In particular taking $\vx=\vy$ shows that $\rT$ is an isometry.
	\item $(\I-\rT^*\rT)^*=\I^*-(\rT^*\rT)^*=\I-\rT^*\rT$ is self-adjoint. By the spectral theorem, there exists an orthonormal basis of $V$ of eigenvectors of $\I-\rT^*\rT$. For any such $\vx$ with (real) eigenvalue $\lambda$,
	\begin{align*}
		0=&\Nm{\vx}^2-\Nm{\rT(\vx)}^2=\ip{\vx,\vx}-\ip{\rT(\vx),\rT(\vx)}=\ip{\vx,(\I-\rT^*\rT)\vx}=\lambda\Nm{\vx}^2\\
		&\implies\lambda=0
	\end{align*}
	Since $\I-\rT^*\rT=0$ on a basis, $\rT^*\rT=\I$. Since $V$ is finite-dimensional, we also have $\rT\rT^*=\I$ whence $\rT$ is unitary.\qedhere
\end{enumerate} 
\end{proof}

The finite-dimensional restriction is important in part 2: we use the existence of adjoints, the spectral theorem, and that a left-inverse is also a right-inverse. See Exercise \ref{ex:isometrynotunitary} for an example of a \emph{non-unitary} isometry in infinite dimensions.\smallbreak
The proof shows a little more:

\begin{cor}{}{unitaryextra}
On a finite-dimensional space, being unitary is equivalent to each of the following:
\begin{enumeratea}\itemsep0pt
	\item Preservation of the inner product\footnotemark $(\dag)$. 
 	\item The existence of an orthonormal basis $\beta=\{\vect{\vw}{n}\}$ such that $\rT(\beta)=\{\rT(\vw_1),\ldots,\rT(\vw_n)\}$ is also orthonormal.
 	\item That \emph{every} orthonormal basis $\beta$ of $V$ is mapped to an orthonormal basis $\rT(\beta)$.
\end{enumeratea} 
\end{cor}

\footnotetext{In particular, in a real inner product space isometries also preserve the \emph{angle} $\theta$ between vectors since $\cos\theta=\frac{\ip{\vx,\vy}}{\Nm\vx\Nm\vy}$. %$(\dag)$ is in fact equivalent to being an isometry in infinite dimensions: recall the \emph{polarization identity\ldots}
}

While (a) is simply $(\dag)$, claims (b) and (c) are also worth proving explicitly: see Exercise \ref{ex:orthbasis}. If $\beta$ is the standard orthonormal basis of $\F^n$ and $\rT=\rL_A$, then the columns of $A$ form the orthonormal set $\rT(\beta)$. This makes identifying unitary/orthogonal matrices easy:

%Before seeing some examples, we state a result which makes identifying these matrices very easy.

\begin{cor}{}{}
A matrix $A\in M_n(\R)$ is orthogonal if and only if its columns form an orthonormal basis of $\R^n$ with respect to the standard (dot) inner product.\\[5pt]
A matrix $A\in M_n(\C)$ is unitary if and only if its columns form an orthonormal basis of $\C^n$ with respect to the standard (Hermitian) inner product.
\end{cor}

% \begin{proof}
% Suppose $\beta$ and $T(\beta)$ are orthonormal: then for any $\vw_j,\vw_k\in\beta$, we have
% \[\delta_{jk}=\ip{\vw_j,\vw_k}\quad\text{and}\quad \delta_{jk}=\ip{T(\vw_j),T(\vw_k)}=\ip{T^*T(\vw_j),\vw_k}\]
% Since $\beta$ is a basis,
% \[\ip{\vw_j-T^*T(\vj),\vw_k}=0\implies \forall \vw\in V,\ \ip{\vw_j-T^*T(\vw_j),\vw}=0\implies T^*T(\vw_j)=\vw_j\]
% Since $T^*T=I$ on a basis, it follows that $T^*T=I$ and $T$ is unitary.\\[5pt]
% If $T$ is unitary and $\beta$ is any orthonormal basis, $(\dag)$ immediately says that $T(\beta)$ is orthonormal. 
% \end{proof}

\goodbreak

\begin{examples}{}{orthomatrix}
\exstart The matrix $A_\theta=\begin{smatrix}
	\cos\theta&-\sin\theta\\
	\sin\theta&\cos\theta
  \end{smatrix}\in M_2(\R)$ is orthogonal for any $\theta$. Example \ref{ex:easyrotate} is this with $\theta=\tan^{-1}\frac 34=\sin^{-1}\frac 35=\cos^{-1}\frac 45$. More generally (Exercise \ref{ex:rotrefmatrix}), it can be seen that every real orthogonal $2\times 2$ matrix has the form $A_\theta$ or \[B_\theta=\begin{pmatrix}
	\cos\theta&\sin\theta\\
	\sin\theta&-\cos\theta
	\end{pmatrix}\]
	for some angle $\theta$. The effect of the $\rL_{A_\theta}$ is to rotate counter-clockwise by $\theta$, while that of $\rL_{B_\theta}$ is to reflect across the line making angle $\frac 12\theta$ with the positive $x$-axis.

\begin{enumerate}\setcounter{enumi}{1}
	\item $A=\frac 1{\sqrt 6}\begin{smatrix}
	\sqrt 2&\sqrt 3&1\\
	\sqrt 2&0&-2\\
	-\sqrt 2&\sqrt 3&-1
	\end{smatrix}\in M_3(\R)$ is orthogonal: check the columns!.
	
	\item $A=\frac 1{\sqrt 2}\begin{smatrix}1&i\\i&1\end{smatrix}\in M_2(\C)$ is unitary: indeed it maps the standard basis to the orthonormal basis 
  \[\rT(\beta)=\left\{\frac 1{\sqrt 2}\twovec 1i, \frac 1{\sqrt 2}\twovec i1\right\}\]
  It is also easy to check that the characteristic polynomial is
  \[p(t)=\det\begin{pmatrix}
  \frac 1{\sqrt 2}-t&\frac i{\sqrt 2}\\
  \frac i{\sqrt 2}&\frac 1{\sqrt 2}-t
  \end{pmatrix} =\left(t-\frac 1{\sqrt 2}\right)^2+\frac 12\implies t=\frac 1{\sqrt 2}(1\pm i) =e^{\pm\pi i/4}\]
  whence the eigenvalues of $\rT$ both have modulus 1.
%   \item Consider $\beta=\{\cos x,\sin x\}$ and $V=\Span\beta$ with $\ip{f,g}=\frac 1\pi\int_{-\pi}^\pi f(t)g(t)\,\dt$. Then $\beta$ is an orthonormal basis. Define $T\in\cL(V)$ by
% \[T(a\cos x+b\sin x)=-b\cos x+a\sin x\]
% This is an orthogonal operator: indeed, writing $c=\cos x$ and $s=\sin x$, we see that
% \begin{gather*}
% \begin{aligned}
% \ip{T^*(ac+bs),pc+qs}&=\ip{ac+bs,T(pc+qs)}=\ip{ac+bs,-qc+ps}\\
% &=-aq+bp=\ip{bc-as,pc+qs}
% \end{aligned}\\
% \implies T^*(a\cos x+b\sin x)=b\cos x-a\sin x\\
% \implies TT^*(a\cos x+b\sin x)=T(b\cos x-a\sin x)=a\cos x+b\sin x
% \end{gather*}
% so that $TT^*=I$. The orthonormal basis $\beta=\{\cos x,\sin x\}$ is mapped to $T(\beta)=\{-\sin x,\cos x\}$, which is clearly also orthonormal. It is easy to check however, that $T$ has no eigenvalues. This example is really just multiplication by $[T]_\beta=\begin{smatrix}0&-1\\1&0\end{smatrix}$ in disguise!

\item\label{ex:infdimunitary} Here is an example of an infinite-dimensional unitary operator. On the space $C[-\pi,\pi]$, the function $\rT(f(x))=e^{ix}f(x)$ is linear. Moreover
\[\ip{e^{ix}f(x),g(x)}=\frac 1{2\pi}\int_{-\pi}^\pi e^{ix}f(x)\cl{g(x)}\,\dx =\frac 1{2\pi}\int_{-\pi}^\pi f(x)\cl{e^{-ix}g(x)}\,\dx =\ip{f(x),e^{-ix}g(x)}\]
whence $\rT^*(f(x))=e^{-ix}f(x)$. Indeed $\rT^*=\rT^{-1}$ and so $\rT$ is a unitary operator.\smallbreak
Since $C[-\pi,\pi]$ is infinite-dimensional, we don't expect all parts of the Corollary to hold:
\begin{itemize}
  \item[(a)] $\rT$ does preserve the inner product.
  \item[(b), (c)] $C[-\pi,\pi]$ doesn't have an orthonormal basis; there is no orthonormal set $\beta=\{f_k\}$ so that every continuous function is a \emph{finite} linear combination.\footnotemark We cannot therefore claim that $\rT$ maps orthonormal bases to orthonormal bases!
\end{itemize}
\end{enumerate}
\end{examples}

\footnotetext{An \emph{infinite} orthonormal set $\beta=\{f_k:k\in\Z\}$ can be found so that every function $f$ `equals' an \emph{infinite series} in the sense that $\Nm{f-\sum a_kf_k}=0$. Since these are not finite sums, $\beta$ isn't strictly a basis, though it isn't uncommon for it to be so described. Moreover, given that the norm is defined by an integral, this also isn't a claim that $f$ and $\sum a_kf_k$ are equal as functions. Indeed the infinite series need not be continuous! For these reasons, when working with Fourier series, one tends to consider a broader class than the continuous functions.}

% \subsubsection*{Unitary and Orthogonal Matrices}
% 
% Suppose $\beta$ is an orthonormal basis of a finite-dimensional $V$ and that $\rT\in\cL(V)$ is unitary. If $A=[\rT]_\beta$, then
% \[A^*A=I=AA^*\]
% It therefore makes sense to extend Definition \ref{defn:unitary} to matrices.
% 
% \begin{defn}{}{unitarymatrix}
% A \emph{unitary matrix} is a matrix $A\in M_n(\C)$ satisfying $A^*A=I$.\\
% An \emph{orthogonal matrix} satisfies $A^TA=I$.
% \end{defn}
% 
% Real unitary matrices are always referred to as orthogonal.\footnote{Complex orthogonal matrices exist (see Exercise \ref{ex:comporth}), though these do not have the same relationship to complex inner products as do unitary matrices.} Note that we don't need the condition $AA^*=I$ for matrices: $A^*A=I$ forces $A$ to be invertible with $A^{-1}=A^*$. %Both conditions are required only for unitary/orthogonal operators on \emph{infinite-dimensional} spaces.
% 
% Before seeing some examples, we state a result which makes identifying these matrices very easy.
% 
% \begin{cor}{}{}
% A matrix $A\in M_n(\R)$ is orthogonal if and only if its columns form an orthonormal basis of $\R^n$ with respect to the standard (dot) inner product.\\[5pt]
% A matrix $A\in M_n(\C)$ is unitary if and only if its columns form an orthonormal basis of $\C^n$ with respect to the standard (hermitian) inner product.
% \end{cor}
% 
% The result follows immediately from Corollary \ref{cor:unitaryextra}: let $\beta=\{\vect{\ve}{n}\}$ be the standard (orthonormal) basis and $\rT=L_A$, then the columns of $A$ are precisely the orthonormal vectors $\rT(\ve_j)$. We can also check this directly by observing that
% \[A=(\vw_1\cdots\vw_n)\implies A^*A=\threevec{\vw_1^*}{\vdots}{\vw_n^*}(\vw_1\ \cdots \vw_n) \implies (A^*A)_{jk}=\vw_j^*\vw_k=\ip{\vw_k,\vw_j}\]
% Clearly $A^*A=I$ if and only if $\{\vect{\vw}{n}\}$ is orthonormal.



\clearpage

\subsubsection*{Unitary and Orthogonal Equivalence}

Suppose $A\in M_n(\R)$ is symmetric (self-adjoint) $A^T=A$. By the spectral theorem, $A$ has an orthonormal eigenbasis $\beta=\{\vect{\vw}{n}\}$: $A\vw_j=\lambda_j\vw_j$. Arranging the eigenbasis as the columns of a matrix, we see that the columns of $U=(\vw_1\cdots \vw_n)$ are orthonormal and so $U$ is an orthogonal matrix. We can therefore write
\[A=UDU^{-1}=U\begin{pmatrix}
\lambda_1&\cdots&0\\
%0&e^{i\theta_2}&&0\\
\vdots&\ddots&\vdots\\
0&\cdots&\lambda_n
\end{pmatrix}U^T\]
A similar approach works if $A\in M_n(\C)$ is normal: we now have $A=UDU^*$ where $U$ is unitary.

\begin{example}{}{easynormal}
The matrix $A=\begin{smatrix}
1+i&1+i\\
-1-i&1+i
\end{smatrix}$ is normal as can easily be checked. Its characteristic polynomial is
\[p(t)=t^2-2(1+i)t+4i=(t-2i)(t-2)\]
with corresponding orthonormal eigenvectors
\[\vw_2=\frac 1{\sqrt 2}\twovec 1{-i},\qquad \vw_{2i}=\frac 1{\sqrt 2}\twovec 1i\]
We conclude that
\[A=\frac 1{\sqrt 2}\begin{pmatrix}
1&1\\
-i&i
\end{pmatrix}
\begin{pmatrix}
2&0\\
0&2i
\end{pmatrix}
\left[\frac 1{\sqrt 2}\begin{pmatrix}
1&1\\
-i&i
\end{pmatrix}\right]^{-1}
=\begin{pmatrix}
1&1\\
-i&i
\end{pmatrix}
\begin{pmatrix}
1&0\\
0&i
\end{pmatrix}
\begin{pmatrix}
1&i\\
1&-i
\end{pmatrix}\]
\end{example}

This is an example of unitary equivalence.

\begin{defn}{}{}
Square matrices $A,B$ are \emph{unitarily equivalent} if there exists a unitary matrix $U$ such that $B=U^*AU$. \emph{Orthogonal equivalence} is similar: $B=U^TAU$.
\end{defn}

The above discussion proves half the following:

\begin{thm}{}{}
$A\in M_n(\C)$ is normal if and only if it is unitarily equivalent to a diagonal matrix (the matrix of its eigenvalues).\smallbreak
$A\in M_n(\R)$ is self-adjoint (symmetric) if and only if it is orthogonally equivalent to a diagonal matrix. 
\end{thm}

\begin{proof}
We've already observed the $(\Rightarrow)$ direction.\smallbreak
For the converse, let $D$ be diagonal, $U$ unitary, and $A=U^*DU$. Then
\begin{align*}
A^*A&=(U^*DU)^*U^*DU=U^*D^*UU^*DU=U^*\cl DDU=U^*D\cl DU =U^*DUU^*\cl DU\\
&=U^*DU(U^*DU)^* =AA^*
\end{align*}
since $U^*=U^{-1}$ and because diagonal matrices commute: $\cl DD=D\cl D$.\smallbreak
In the special case where $A$ is real and $U$ is orthogonal, then $A$ is symmetric:
\[A^T=(U^TDU)^T=U^TD^TU=U^TDU=A\tag*{\qedhere}\]
\end{proof}

\goodbreak



\begin{exercises}
\exstart\label{ex:orthequiv1} For each matrix $A$ find an orthogonal or unitary $U$ and a diagonal $D=U^*AU$.

\begin{enumerate}\setcounter{enumi}{1}  
  \item[]\begin{enumerate}
    \item $\begin{smatrix}
    1&2\\2&1
    \end{smatrix}$
    \qquad
    (b)\ \ $\begin{smatrix}
    0&-1\\1&0
    \end{smatrix}$
    \qquad
    (c)\ \ $\begin{smatrix}
    2&3-3i\\3+3i&5
    \end{smatrix}$
    \qquad
    (d)\ \ $\begin{smatrix}
    2&1&1\\1&2&1\\1&1&2
    \end{smatrix}$
  \end{enumerate}
  
  \item Which of the following pairs are unitarily/orthogonally equivalent? Explain your answers.
  \begin{enumerate}
    \item $A=\begin{smatrix}
    0&1\\1&0
    \end{smatrix}$ and $B=\begin{smatrix}
    0&2\\2&0
    \end{smatrix}$\qquad\quad (b)\ \ $A=\begin{smatrix}
    0&-1&0\\1&0&0\\0&0&1
    \end{smatrix}$ and $B=\begin{smatrix}
    2&0&0\\0&-1&0\\0&0&0
    \end{smatrix}$
    \item[(c)] $A=\begin{smatrix}
    0&-1&0\\1&0&0\\0&0&1
    \end{smatrix}$ and $B=\begin{smatrix}
    1&0&0\\0&i&0\\0&0&-i
    \end{smatrix}$
  \end{enumerate}
	
	\item Let $a,b\in\C$ be such that $\nm a^2+\nm b^2=1$. Prove that every $2\times 2$ matrix of the form $\begin{smatrix}
	a&-e^{i\theta}\cl b\\
	b&e^{i\theta}\cl a
	\end{smatrix}$ is unitary. Are these all the unitary $2\times 2$ matrices? Prove or disprove.
	
	
	\item If $A,B$ are orthogonal/unitary, prove that $AB$ and $A^{-1}$ are also orthogonal/unitary.\\[2pt]
	(\emph{This proves that orthogonal/unitary matrices are groups under matrix multiplication})
  
  \item\label{ex:comporth} Check that $A=\frac 13\begin{smatrix}
	5&-4i\\
	4i&5
  \end{smatrix}\in M_2(\C)$ satisfies $A^TA=I$ (it is a \emph{complex orthogonal} matrix).\\[2pt]
  (\emph{These don't have the same nice relationship with inner products, and are thus less useful to us}) 
  
  \item\label{ex:rotrefmatrix} Supply the details of Exercise \hyperref[ex:orthomatrix]{\ref*{ex:orthomatrix}.1}.\\[2pt]
  (\emph{Hints: $\beta=\{\vi,\vj\}$ is orthonormal, whence $\{A\vi,A\vj\}$ must be orthonormal. Now draw pictures to compute the result of rotating and reflecting the vectors $\vi$ and $\vj$.})
  
	
	\item Show that the linear map in Example \hyperref[ex:infdimunitary]{\ref*{ex:orthomatrix}.\ref*{ex:infdimunitary}} has no eigenvectors.
	
% 	However, in contrast to a unitary operator on a finite-dimensional complex space, $\rT$ has no eigenvalues/eigenvectors, since
% 	\[\rT(f)=\lambda f\iff \forall x,\ e^{ix}f(x)=\lambda f(x) \iff f(x)\equiv 0\]
  
  \item\label{ex:unitarymatrix} Prove that $A\in M_n(\C)$ has an orthonormal basis of eigenvectors whose eigenvalues have modulus 1, if and only if $A$ is unitary.
  
  
  \item\label{ex:orthbasis} Prove parts (b) and (c) of Corollary \ref{cor:unitaryextra} for a finite-dimensional inner product space:
  \begin{enumerate}
    \item If $\beta$ is an orthonormal basis such that $\rT(\beta)$ is orthonormal, then $\rT$ is unitary.
    \item If $\rT$ is unitary, and $\eta$ is an orthonormal basis, then $\rT(\eta)$ is an orthonormal basis.
  \end{enumerate}
  
  \item Let $\rT$ be a linear operator on a finite-dimensional inner product space $V$. If $\Nm{\rT(\vx)}=\Nm{\vx}$ for all $\vx$ in some orthonormal basis of $V$, must $\rT$ be unitary? Prove or disprove.
  
  \item Let $\rT$ be a unitary operator on an inner product space $V$ and let $W$ be a finite-dimensional $\rT$-invariant subspace of $V$. Prove:
  \begin{enumerate}
    \item $\rT(W)=W$\quad (\emph{Hint: show that $\at\rT W$ is injective});
    \item $W^\perp$ is $\rT$-invariant.
	\end{enumerate}
	
	\item Let $W$ a subspace of an inner product space $V$ such that $V=W\oplus W^\perp$. Define $\rT\in\cL(V)$ by $\rT(\vu+\vw)=\vu-\vw$ where $\vu\in W$ and $\vw\in W^\perp$. Prove that $\rT$ is unitary and self-adjoint.
	
	\item\label{ex:isometrynotunitary} In the inner product space $\ell^2$ of square-summable sequences, consider the linear operator
	$\rT(x_1,x_2,\ldots)=(0,x_1,x_2,\ldots)$. Prove that $\rT$ is an isometry and compute its adjoint. Check that $\rT$ is non-invertible and non-unitary.
	
	\item Prove Schur's Lemma for matrices. Every $A\in M_n(\R)$ is orthogonally equivalent and every $A\in M_n(\C)$ is unitarily equivalent to an upper triangular matrix. 
\end{enumerate}
\end{exercises}